% =====================================================================
%  Huffman frente a LZW: comparación teórica y experimental
%  Universidad de Lima — Trabajo de Investigación 2026-2
%  Compilar con: ./build.sh  (tectonic ejecuta BibTeX automáticamente)
% =====================================================================
\documentclass[11pt,a4paper]{article}

% ---------------------------------------------------------------- idioma y tipografía
\usepackage{iftex}
\ifXeTeX
  % Tectonic compila con XeTeX: Latin Modern en OpenType mediante fontspec, para que
  % los caracteres UTF-8 (por ejemplo, ¿ y ¡) se rendericen correctamente.
  \usepackage{fontspec}
\else
  \usepackage[T1]{fontenc}
  \usepackage[utf8]{inputenc}
  \usepackage{lmodern}
\fi
\usepackage[english,spanish,es-noshorthands,es-tabla,es-nodecimaldot]{babel}
\usepackage{microtype}
\usepackage[margin=2.5cm,headheight=14pt]{geometry}

% ---------------------------------------------------------------- matemática, unidades, tablas
\usepackage{amsmath}
\usepackage{amssymb}
\usepackage{siunitx}
\sisetup{group-separator={\,}, group-minimum-digits=4}
\usepackage{booktabs}
\usepackage{array}
\usepackage{multirow}
\usepackage{enumitem}

% ---------------------------------------------------------------- gráficos
\usepackage{graphicx}
\usepackage[dvipsnames]{xcolor}
\usepackage{tikz}
\usetikzlibrary{positioning,arrows.meta,shapes.geometric,calc,fit,babel}
\usepackage{pgfplots}
\usepackage{pgfplotstable}
\pgfplotsset{compat=1.18}
\usepackage{caption}
\usepackage{subcaption}
\captionsetup{font=small,labelfont=bf,format=hang}

% ---------------------------------------------------------------- algoritmos y código
\usepackage{float}
\usepackage{algorithm}
\usepackage{algpseudocode}
\usepackage{listings}

% ---------------------------------------------------------------- cajas de aporte
\usepackage[most]{tcolorbox}

% ---------------------------------------------------------------- encabezados
\usepackage{fancyhdr}

% ---------------------------------------------------------------- enlaces (al final)
\usepackage[colorlinks=true,linkcolor=blue!55!black,citecolor=green!40!black,
            urlcolor=blue!60!black,bookmarksnumbered=true]{hyperref}
\usepackage[spanish,nameinlink]{cleveref}

% =====================================================================
%  INTEGRANTES Y CÓDIGO DE COLORES — ÚNICO LUGAR PARA EDITAR LOS NOMBRES
% =====================================================================
\newcommand{\integranteA}{Integrante A (Apellido, Nombre)}
\newcommand{\integranteB}{Integrante B (Apellido, Nombre)}
\newcommand{\integranteC}{Integrante C (Apellido, Nombre)}
\newcommand{\integranteD}{Integrante D (Apellido, Nombre)}
\newcommand{\docente}{Docente del curso (Apellido, Nombre)}

\colorlet{memberA}{blue!70!black}
\colorlet{memberB}{green!45!black}
\colorlet{memberC}{orange!75!black}
\colorlet{memberD}{violet!75!black}

% Entorno de autoría: texto en el color del integrante y barra lateral fina.
% Uso: \begin{aporte}{A} ... \end{aporte}.  Las figuras, tablas y algoritmos
% dentro de un aporte deben usar la posición [H] (una caja no admite flotantes).
\newtcolorbox{aporte}[1]{%
  enhanced jigsaw, breakable, parbox=false,
  colback=white, coltext=member#1,
  frame hidden, boxrule=0pt, sharp corners,
  borderline west={1.2pt}{0pt}{member#1},
  left=9pt, right=0pt, top=1pt, bottom=1pt,
  before skip=4pt, after skip=8pt,
  before upper={\noindent{\scriptsize\sffamily\bfseries
    Aporte: \csname integrante#1\endcsname}\par\smallskip}}

% Muestra de color para la tabla de distribución.
\newcommand{\muestra}[1]{\textcolor{member#1}{\rule{1.6em}{0.9em}}}

% =====================================================================
%  RESULTADOS GENERADOS (results/ se copia desde ../results en build.sh)
% =====================================================================
\newcommand{\resmissing}{\textcolor{red}{\textbf{??}}}
\IfFileExists{results/macros.tex}{% Generated by BenchmarkRunner; do not edit.
\newcommand{\resJvm}{OpenJDK 64-Bit Server VM 21.0.13}
\newcommand{\resCpu}{Intel(R) Core(TM) Ultra 7 265K}
\newcommand{\resCores}{20}
\newcommand{\resHeap}{2944}
\newcommand{\resDate}{2026-11-09}
\newcommand{\resSlopeUniformHuffEnc}{1.03}
\newcommand{\resRtwoUniformHuffEnc}{0.993}
\newcommand{\resNspbUniformHuffEnc}{5.37}
\newcommand{\resSlopeUniformHuffDec}{1.04}
\newcommand{\resRtwoUniformHuffDec}{0.998}
\newcommand{\resNspbUniformHuffDec}{9.14}
\newcommand{\resSlopeUniformLzwEnc}{1.05}
\newcommand{\resRtwoUniformLzwEnc}{0.997}
\newcommand{\resNspbUniformLzwEnc}{13.05}
\newcommand{\resSlopeUniformLzwDec}{1.01}
\newcommand{\resRtwoUniformLzwDec}{0.997}
\newcommand{\resNspbUniformLzwDec}{5.11}
\newcommand{\resSlopeSkewedHuffEnc}{1.05}
\newcommand{\resRtwoSkewedHuffEnc}{0.996}
\newcommand{\resNspbSkewedHuffEnc}{4.50}
\newcommand{\resSlopeSkewedHuffDec}{1.04}
\newcommand{\resRtwoSkewedHuffDec}{0.997}
\newcommand{\resNspbSkewedHuffDec}{11.17}
\newcommand{\resSlopeSkewedLzwEnc}{1.08}
\newcommand{\resRtwoSkewedLzwEnc}{0.987}
\newcommand{\resNspbSkewedLzwEnc}{10.54}
\newcommand{\resSlopeSkewedLzwDec}{1.01}
\newcommand{\resRtwoSkewedLzwDec}{0.990}
\newcommand{\resNspbSkewedLzwDec}{3.35}
\newcommand{\resSlopeSynthtextHuffEnc}{1.07}
\newcommand{\resRtwoSynthtextHuffEnc}{0.998}
\newcommand{\resNspbSynthtextHuffEnc}{6.26}
\newcommand{\resSlopeSynthtextHuffDec}{1.04}
\newcommand{\resRtwoSynthtextHuffDec}{0.998}
\newcommand{\resNspbSynthtextHuffDec}{10.64}
\newcommand{\resSlopeSynthtextLzwEnc}{1.08}
\newcommand{\resRtwoSynthtextLzwEnc}{0.998}
\newcommand{\resNspbSynthtextLzwEnc}{14.72}
\newcommand{\resSlopeSynthtextLzwDec}{1.00}
\newcommand{\resRtwoSynthtextLzwDec}{0.991}
\newcommand{\resNspbSynthtextLzwDec}{4.12}
\newcommand{\resSlopeDickensHuffEnc}{0.94}
\newcommand{\resRtwoDickensHuffEnc}{0.994}
\newcommand{\resNspbDickensHuffEnc}{4.53}
\newcommand{\resSlopeDickensHuffDec}{0.95}
\newcommand{\resRtwoDickensHuffDec}{0.997}
\newcommand{\resNspbDickensHuffDec}{9.59}
\newcommand{\resSlopeDickensLzwEnc}{0.88}
\newcommand{\resRtwoDickensLzwEnc}{0.985}
\newcommand{\resNspbDickensLzwEnc}{10.34}
\newcommand{\resSlopeDickensLzwDec}{0.85}
\newcommand{\resRtwoDickensLzwDec}{0.991}
\newcommand{\resNspbDickensLzwDec}{3.17}
\newcommand{\resEntropyHuffBpbKTwo}{1.00}
\newcommand{\resEntropyLzwBpbKTwo}{1.13}
\newcommand{\resEntropyHuffEncMsKTwo}{26.95}
\newcommand{\resEntropyHuffDecMsKTwo}{8.60}
\newcommand{\resEntropyLzwEncMsKTwo}{60.60}
\newcommand{\resEntropyLzwDecMsKTwo}{13.17}
\newcommand{\resEntropyHuffBpbKFour}{2.00}
\newcommand{\resEntropyLzwBpbKFour}{2.22}
\newcommand{\resEntropyHuffEncMsKFour}{24.33}
\newcommand{\resEntropyHuffDecMsKFour}{15.12}
\newcommand{\resEntropyLzwEncMsKFour}{55.63}
\newcommand{\resEntropyLzwDecMsKFour}{15.34}
\newcommand{\resEntropyHuffBpbKEight}{3.00}
\newcommand{\resEntropyLzwBpbKEight}{3.37}
\newcommand{\resEntropyHuffEncMsKEight}{23.71}
\newcommand{\resEntropyHuffDecMsKEight}{21.87}
\newcommand{\resEntropyLzwEncMsKEight}{58.16}
\newcommand{\resEntropyLzwDecMsKEight}{18.45}
\newcommand{\resEntropyHuffBpbKSixteen}{4.00}
\newcommand{\resEntropyLzwBpbKSixteen}{4.61}
\newcommand{\resEntropyHuffEncMsKSixteen}{25.35}
\newcommand{\resEntropyHuffDecMsKSixteen}{28.96}
\newcommand{\resEntropyLzwEncMsKSixteen}{78.62}
\newcommand{\resEntropyLzwDecMsKSixteen}{24.02}
\newcommand{\resEntropyHuffBpbKThirtyTwo}{5.00}
\newcommand{\resEntropyLzwBpbKThirtyTwo}{5.90}
\newcommand{\resEntropyHuffEncMsKThirtyTwo}{28.07}
\newcommand{\resEntropyHuffDecMsKThirtyTwo}{37.34}
\newcommand{\resEntropyLzwEncMsKThirtyTwo}{80.49}
\newcommand{\resEntropyLzwDecMsKThirtyTwo}{27.25}
\newcommand{\resEntropyHuffBpbKSixtyFour}{6.00}
\newcommand{\resEntropyLzwBpbKSixtyFour}{7.37}
\newcommand{\resEntropyHuffEncMsKSixtyFour}{30.39}
\newcommand{\resEntropyHuffDecMsKSixtyFour}{43.70}
\newcommand{\resEntropyLzwEncMsKSixtyFour}{85.42}
\newcommand{\resEntropyLzwDecMsKSixtyFour}{26.27}
\newcommand{\resEntropyHuffBpbKOneTwentyEight}{7.00}
\newcommand{\resEntropyLzwBpbKOneTwentyEight}{8.53}
\newcommand{\resEntropyHuffEncMsKOneTwentyEight}{29.95}
\newcommand{\resEntropyHuffDecMsKOneTwentyEight}{49.94}
\newcommand{\resEntropyLzwEncMsKOneTwentyEight}{85.42}
\newcommand{\resEntropyLzwDecMsKOneTwentyEight}{32.12}
\newcommand{\resEntropyHuffBpbKTwoFiftySix}{8.00}
\newcommand{\resEntropyLzwBpbKTwoFiftySix}{11.00}
\newcommand{\resEntropyHuffEncMsKTwoFiftySix}{28.18}
\newcommand{\resEntropyHuffDecMsKTwoFiftySix}{57.26}
\newcommand{\resEntropyLzwEncMsKTwoFiftySix}{107.11}
\newcommand{\resEntropyLzwDecMsKTwoFiftySix}{42.08}
\newcommand{\resCorpusFiles}{26}
\newcommand{\resLzwWins}{24}
\newcommand{\resCorpusHuffBpb}{5.12}
\newcommand{\resCorpusLzwBpb}{3.30}
\newcommand{\resCorpusDeflateBpb}{2.56}
\newcommand{\resLzwBitsBestDickens}{16}
\newcommand{\resLzwBitsBestXml}{16}
\newcommand{\resLzwBitsBestEcoli}{16}
}{}
% Valores de respaldo: solo se usan si el benchmark aún no generó la macro.
\providecommand{\resJvm}{\resmissing}
\providecommand{\resCpu}{\resmissing}
\providecommand{\resCores}{\resmissing}
\providecommand{\resHeap}{\resmissing}
\providecommand{\resDate}{\resmissing}
\providecommand{\resCorpusFiles}{\resmissing}
\providecommand{\resLzwWins}{\resmissing}
\providecommand{\resCorpusHuffBpb}{\resmissing}
\providecommand{\resCorpusLzwBpb}{\resmissing}
\providecommand{\resCorpusDeflateBpb}{\resmissing}
\providecommand{\resLzwBitsBestDickens}{\resmissing}
\providecommand{\resLzwBitsBestXml}{\resmissing}
\providecommand{\resLzwBitsBestEcoli}{\resmissing}
\providecommand{\resSlopeDickensHuffEnc}{\resmissing}
\providecommand{\resSlopeDickensHuffDec}{\resmissing}
\providecommand{\resSlopeDickensLzwEnc}{\resmissing}
\providecommand{\resSlopeDickensLzwDec}{\resmissing}
\providecommand{\resRtwoDickensHuffEnc}{\resmissing}
\providecommand{\resRtwoDickensHuffDec}{\resmissing}
\providecommand{\resRtwoDickensLzwEnc}{\resmissing}
\providecommand{\resRtwoDickensLzwDec}{\resmissing}
\providecommand{\resNspbDickensHuffEnc}{\resmissing}
\providecommand{\resNspbDickensHuffDec}{\resmissing}
\providecommand{\resNspbDickensLzwEnc}{\resmissing}
\providecommand{\resNspbDickensLzwDec}{\resmissing}
\providecommand{\resEntropyHuffBpbKTwo}{\resmissing}
\providecommand{\resEntropyHuffBpbKTwoFiftySix}{\resmissing}
\providecommand{\resEntropyLzwBpbKTwo}{\resmissing}
\providecommand{\resEntropyLzwBpbKTwoFiftySix}{\resmissing}
\providecommand{\resEntropyHuffDecMsKTwo}{\resmissing}
\providecommand{\resEntropyHuffDecMsKTwoFiftySix}{\resmissing}
\providecommand{\resEntropyLzwDecMsKTwo}{\resmissing}
\providecommand{\resEntropyLzwDecMsKTwoFiftySix}{\resmissing}

% Inserta #2 solo si existe results/#1; si no, deja un aviso visible.
\newcommand{\siresultado}[2]{%
  \IfFileExists{results/#1}{#2}{\todoresults{Falta el archivo \texttt{results/#1}; ejecutar \texttt{ant benchmark}.}}}

% Marcador de interpretación pendiente (lo completa el equipo tras la corrida).
\newcommand{\todoresults}[1]{%
  \par\smallskip\noindent
  \fcolorbox{red}{red!4}{\parbox{\dimexpr\linewidth-2\fboxsep-2\fboxrule\relax}{%
    \color{red!80!black}\small\textbf{Pendiente (resultados):} #1}}\par\smallskip}

% =====================================================================
%  ESTILOS
% =====================================================================
\pgfplotsset{
  informe/.style={
    width=\linewidth, height=0.72\linewidth,
    grid=major, grid style={gray!25},
    tick label style={font=\scriptsize}, label style={font=\small},
    legend style={font=\scriptsize, draw=gray!50, fill=white, fill opacity=0.9,
                  text opacity=1, cells={anchor=west}},
    legend pos=north west,
    error bars/error bar style={line width=0.4pt},
    error bars/error mark options={rotate=90, mark size=1.5pt},
  },
  huffenc/.style={blue!70!black, mark=*, mark size=1.3pt, line width=0.7pt},
  huffdec/.style={blue!70!black, mark=o, mark size=1.5pt, line width=0.7pt, densely dashed},
  lzwenc/.style={red!70!black, mark=square*, mark size=1.2pt, line width=0.7pt},
  lzwdec/.style={red!70!black, mark=square, mark size=1.4pt, line width=0.7pt, densely dashed},
  pendiente/.style={black!60, dashed, line width=0.6pt, mark=none},
}

\definecolor{codebg}{gray}{0.97}
\lstdefinestyle{java}{
  language=Java, basicstyle=\ttfamily\footnotesize, keywordstyle=\color{blue!60!black},
  commentstyle=\color{gray!70!black}\itshape, stringstyle=\color{green!40!black},
  numbers=left, numberstyle=\tiny\color{gray}, numbersep=6pt, stepnumber=1,
  backgroundcolor=\color{codebg}, frame=single, rulecolor=\color{gray!40},
  breaklines=true, showstringspaces=false, tabsize=4, columns=fullflexible,
  xleftmargin=1.5em, framexleftmargin=1.5em, captionpos=b,
}
\lstset{style=java}
\renewcommand{\lstlistingname}{Listado}
\renewcommand{\lstlistlistingname}{Índice de listados}

% Pseudocódigo en español.
\floatname{algorithm}{Algoritmo}
\renewcommand{\listalgorithmname}{Índice de algoritmos}
\algrenewcommand\algorithmicrequire{\textbf{Entrada:}}
\algrenewcommand\algorithmicensure{\textbf{Salida:}}
\algrenewcommand\algorithmicfunction{\textbf{función}}
\algrenewcommand\algorithmicprocedure{\textbf{procedimiento}}
\algrenewcommand\algorithmicend{\textbf{fin}}
\algrenewcommand\algorithmicif{\textbf{si}}
\algrenewcommand\algorithmicthen{\textbf{entonces}}
\algrenewcommand\algorithmicelse{\textbf{si no}}
\algrenewcommand\algorithmicfor{\textbf{para}}
\algrenewcommand\algorithmicforall{\textbf{para cada}}
\algrenewcommand\algorithmicwhile{\textbf{mientras}}
\algrenewcommand\algorithmicdo{\textbf{hacer}}
\algrenewcommand\algorithmicreturn{\textbf{devolver}}

\crefname{table}{tabla}{tablas}
\Crefname{table}{Tabla}{Tablas}
\crefname{figure}{figura}{figuras}
\Crefname{figure}{Figura}{Figuras}
\crefname{algorithm}{algoritmo}{algoritmos}
\Crefname{algorithm}{Algoritmo}{Algoritmos}
\crefname{lstlisting}{listado}{listados}
\Crefname{lstlisting}{Listado}{Listados}
\crefname{section}{sección}{secciones}
\Crefname{section}{Sección}{Secciones}
\crefname{subsection}{sección}{secciones}
\Crefname{subsection}{Sección}{Secciones}
\crefname{appendix}{anexo}{anexos}
\Crefname{appendix}{Anexo}{Anexos}
\crefname{equation}{ecuación}{ecuaciones}
\Crefname{equation}{Ecuación}{Ecuaciones}

\newcommand{\Lbar}{\bar{L}}

\pagestyle{fancy}
\fancyhf{}
\fancyhead[L]{\small\itshape Huffman frente a LZW}
\fancyhead[R]{\small\itshape Universidad de Lima}
\fancyfoot[C]{\small\thepage}
\renewcommand{\headrulewidth}{0.4pt}

\setlist{itemsep=2pt, topsep=4pt}

\hypersetup{
  pdftitle={Huffman frente a LZW: comparación teórica y experimental de algoritmos de compresión sin pérdida},
  pdfauthor={Universidad de Lima, Trabajo de Investigación 2026-2},
  pdfkeywords={compresión sin pérdida, codificación de Huffman, LZW, análisis de complejidad, algoritmia experimental, entropía},
}

% =====================================================================
\begin{document}

% ---------------------------------------------------------------- portada
\begin{titlepage}
  \centering
  \includegraphics[width=4.2cm]{figures/ulima-logo.jpg}\par\vspace{1.0cm}
  {\large\scshape Universidad de Lima\par}\vspace{0.2cm}
  {\normalsize Facultad de Ingeniería — Carrera de Ingeniería de Sistemas\par}\vspace{1.4cm}
  {\normalsize Trabajo de Investigación — Curso: (Área de Ingeniería de Software)\par}\vspace{1.2cm}
  \rule{\linewidth}{0.8pt}\par\vspace{0.4cm}
  {\LARGE\bfseries Huffman frente a LZW:\par}\vspace{0.2cm}
  {\Large comparación teórica y experimental de algoritmos\\ de compresión sin pérdida\par}
  \vspace{0.4cm}\rule{\linewidth}{0.8pt}\par\vspace{1.6cm}
  {\normalsize\bfseries Integrantes\par}\vspace{0.3cm}
  \begin{tabular}{@{}c@{\quad}l@{}}
    \muestra{A} & \textcolor{memberA}{\integranteA}\\[2pt]
    \muestra{B} & \textcolor{memberB}{\integranteB}\\[2pt]
    \muestra{C} & \textcolor{memberC}{\integranteC}\\[2pt]
    \muestra{D} & \textcolor{memberD}{\integranteD}\\
  \end{tabular}\par\vspace{1.2cm}
  {\normalsize\bfseries Docente\par}\vspace{0.2cm}
  {\normalsize \docente\par}
  \vfill
  {\normalsize Lima, Perú — 2026-2\par}
\end{titlepage}

\pagenumbering{roman}

% ---------------------------------------------------------------- resumen
\section*{Resumen}
\addcontentsline{toc}{section}{Resumen}
\begin{aporte}{D}
La compresión sin pérdida reduce el costo de almacenar y transmitir datos sin alterar un
solo bit de la información original. Este trabajo compara dos algoritmos clásicos que
explotan redundancias de naturaleza distinta: la codificación de Huffman, un código prefijo
estático que aprovecha la distribución de frecuencias de los símbolos, y el algoritmo
Lempel--Ziv--Welch (LZW), un método adaptativo de diccionario que reemplaza frases
repetidas por códigos. Ambos se implementaron en Java (proyecto NetBeans con Ant), con
codificación canónica de Huffman y LZW de ancho variable de 9 a 16 bits con reinicio del
diccionario, validados con 43 pruebas JUnit. Se derivó la complejidad temporal y espacial de
cada algoritmo y se contrastó con cuatro experimentos: escalamiento del tiempo con el tamaño
de la entrada (ajuste log--log), barrido de entropía con alfabetos uniformes, evaluación sobre
los corpus Canterbury y Silesia con DEFLATE como referencia, y efecto del tamaño máximo del
diccionario de LZW. Los tiempos se reportan como medianas con intervalos de confianza del
95\,\% tras el calentamiento del compilador JIT.
% TODO-RESULTS: resumir en dos o tres oraciones los hallazgos principales (pendientes, bpb agregados del corpus, ganador por tipo de dato).
\todoresults{Resumir los hallazgos principales: pendientes log--log obtenidas, bits por byte agregados en el corpus y tipo de dato en que gana cada algoritmo.}

\medskip\noindent\textbf{Palabras clave:} compresión sin pérdida, codificación de Huffman,
LZW, análisis de complejidad, algoritmia experimental, entropía.
\end{aporte}

\begin{otherlanguage}{english}
\section*{Abstract}
\addcontentsline{toc}{section}{Abstract}
\begin{aporte}{D}
Lossless compression reduces the cost of storing and transmitting data without altering a
single bit of the original information. This work compares two classical algorithms that
exploit different kinds of redundancy: Huffman coding, a static prefix code that exploits
the symbol frequency distribution, and the Lempel--Ziv--Welch (LZW) algorithm, an adaptive
dictionary method that replaces repeated phrases with codes. Both were implemented in Java
(a NetBeans project built with Ant), using canonical Huffman codes and variable-width LZW
codes from 9 to 16 bits with dictionary reset, and validated with 43 JUnit tests. The time and
space complexity of each algorithm was derived and contrasted through four experiments:
time scaling with input size (log--log fit), an entropy sweep over uniform alphabets, an
evaluation on the Canterbury and Silesia corpora with DEFLATE as a reference, and the effect
of the maximum LZW dictionary size. Times are reported as medians with 95\,\% confidence
intervals after JIT warm-up.
% TODO-RESULTS: summarize the main findings in two or three sentences (English).
\todoresults{Summarize the main findings in English, consistent with the Spanish abstract.}

\medskip\noindent\textbf{Keywords:} lossless compression, Huffman coding, LZW, complexity
analysis, experimental algorithmics, entropy.
\end{aporte}
\end{otherlanguage}

\clearpage
\tableofcontents
\clearpage
\listoffigures
\listoftables
\listofalgorithms
\lstlistoflistings
\clearpage

% ---------------------------------------------------------------- distribución del trabajo
\section*{Distribución del trabajo y código de colores}
\addcontentsline{toc}{section}{Distribución del trabajo y código de colores}
Cada integrante es responsable de una parte del estudio. En el cuerpo del informe, el
contenido de cada sección o subsección aparece en el color de su autor y con una barra
lateral del mismo color, encabezada por el nombre del integrante. El detalle de cada
implementación individual se presenta en el \cref{app:aportes}.

\begin{table}[H]
  \centering
  \caption{Integrantes, color asignado y responsabilidades.}
  \label{tab:distribucion}
  \small
  \begin{tabular}{@{}c l p{9.2cm}@{}}
    \toprule
    Color & Integrante & Responsabilidades \\
    \midrule
    \muestra{A} & \textcolor{memberA}{\integranteA} &
      Codificación de Huffman: marco teórico, análisis de complejidad, implementación
      (\texttt{HuffmanCodec}, \texttt{BitWriter}, \texttt{BitReader}) y sus pruebas;
      discusión de H2 y H4. \\[3pt]
    \muestra{B} & \textcolor{memberB}{\integranteB} &
      Algoritmo LZW: marco teórico, análisis de complejidad, implementación
      (\texttt{LzwCodec}) y sus pruebas; discusión de H3. \\[3pt]
    \muestra{C} & \textcolor{memberC}{\integranteC} &
      Metodología experimental, comparación teórica e hipótesis, arnés de medición
      (\texttt{BenchmarkRunner}, \texttt{Stats}, \texttt{DataGenerator}), interfaz de
      línea de comandos, figuras y tablas de resultados; discusión de H1. \\[3pt]
    \muestra{D} & \textcolor{memberD}{\integranteD} &
      Resumen, introducción, definiciones preliminares, estado del arte, síntesis de las
      preguntas de investigación, conclusiones y bibliografía. \\
    \bottomrule
  \end{tabular}
\end{table}

\clearpage
\pagenumbering{arabic}

% ---------------------------------------------------------------------
%  1. Introducción
% ---------------------------------------------------------------------
\section{Introducción}
\label{sec:introduccion}

\subsection{Contexto}
El volumen de datos que se almacena y transmite crece de forma sostenida, y con él el costo
de discos, memoria y ancho de banda. La compresión de datos reduce ese costo representando la
misma información con menos bits. Cuando la reconstrucción debe ser exacta ---texto, código
fuente, ejecutables, registros científicos o bases de datos---, solo son aceptables los
métodos \emph{sin pérdida}, en los que la descompresión devuelve bit a bit la entrada
original~\cite{sayood2017,salomon2010}. La teoría de la información establece el límite de lo
alcanzable: ningún código puede representar una fuente con menos bits por símbolo que su
entropía~\cite{shannon1948}.

Dos algoritmos ocupan un lugar central en la historia y en la práctica de la compresión sin
pérdida. La codificación de Huffman~\cite{huffman1952} construye un código prefijo óptimo a
partir de las frecuencias de los símbolos, y el algoritmo Lempel--Ziv--Welch
(LZW)~\cite{welch1984}, derivado de LZ78~\cite{ziv1978}, sustituye frases repetidas por
referencias a un diccionario que se construye durante la lectura. Ambos siguen vigentes: la
codificación de Huffman forma parte de DEFLATE~\cite{rfc1951} y de numerosos formatos, y LZW
se emplea en formatos de imagen y de archivo, y es una referencia habitual en dispositivos con
recursos limitados~\cite{jambek2014,bjaras2019}. Comparar ambos enfoques permite entender dos
formas complementarias de explotar la redundancia: la estadística de los símbolos aislados y la
repetición de secuencias~\cite{bell1989,jayasankar2021}.

\subsection{Problema}
Los estudios comparativos de Huffman y LZW disponibles suelen limitarse a un único dominio de
datos o a archivos de pocos tamaños, reportan tiempos de una sola ejecución y no contrastan
las mediciones con la complejidad teórica de los
algoritmos~\cite{jambek2014,fauzan2022,btoush2008}. En consecuencia, no queda claro en qué
condiciones cada algoritmo comprime mejor, si sus tiempos crecen realmente de forma lineal con
el tamaño de la entrada, ni cómo influyen la entropía de la fuente o el tamaño del diccionario.
El problema que aborda este trabajo es, por tanto: \emph{¿cómo se comparan Huffman y LZW, en
tiempo de ejecución y en tasa de compresión, sobre una gran cantidad de datos de naturaleza
diversa, y en qué medida el comportamiento observado coincide con el análisis teórico?}

\subsection{Objetivos}
\paragraph{Objetivo general.} Comparar teórica y experimentalmente la codificación de Huffman y
el algoritmo LZW en tiempo de compresión, tiempo de descompresión y tasa de compresión,
contrastando las mediciones con el análisis de complejidad de cada algoritmo.

\paragraph{Objetivos específicos.}
\begin{enumerate}[label=\textbf{OE\arabic*.}, leftmargin=*]
  \item Implementar en Java ambos algoritmos (Huffman canónico y LZW de ancho variable de 9 a
        16 bits) y verificar su corrección con una batería de pruebas unitarias que incluya
        ida y vuelta (\emph{round trip}) sin pérdida en el 100\,\% de los casos.
  \item Derivar la complejidad temporal y espacial de codificación y decodificación de cada
        algoritmo y formular hipótesis contrastables a partir de ella.
  \item Medir el tiempo de ejecución en función del tamaño de la entrada desde
        16~KiB ($2^{14}$ bytes) hasta 16~MiB ($2^{24}$ bytes) en cuatro escenarios de datos (hasta 8~MiB en el texto real) y estimar
        el exponente de crecimiento mediante regresión log--log con su coeficiente $R^2$.
  \item Medir la tasa de compresión (bits por byte) frente a la entropía de orden cero en un
        barrido controlado de alfabetos y en los corpus estándar Canterbury y Silesia.
  \item Cuantificar el efecto del tamaño máximo del diccionario de LZW (de 9 a 16 bits) sobre la
        tasa de compresión y el tiempo.
\end{enumerate}

\subsection{Preguntas de investigación}
\begin{enumerate}[label=\textbf{RQ\arabic*.}, leftmargin=*]
  \item ¿El tiempo de codificación y decodificación de Huffman y de LZW crece con el tamaño de
        la entrada según lo que predice su análisis de complejidad?
  \item ¿Cómo se relaciona la tasa de compresión de cada algoritmo con la entropía de orden cero
        y con el tipo de dato (fuentes sin memoria frente a datos con frases repetidas)?
  \item ¿Qué efecto tiene el tamaño máximo del diccionario de LZW sobre la tasa de compresión y
        el tiempo de ejecución?
\end{enumerate}

\subsection{Contribuciones}
\begin{itemize}
  \item Dos implementaciones documentadas y probadas con una interfaz común
        (\texttt{Codec}): Huffman canónico estático y LZW de ancho variable con reinicio del
        diccionario.
  \item Un análisis de complejidad de ambos algoritmos ligado a las decisiones concretas de
        implementación (montículo binario, decodificación canónica, tabla hash de
        direccionamiento abierto).
  \item Un protocolo experimental reproducible, con calentamiento del JIT, repeticiones,
        medianas e intervalos de confianza~\cite{georges2007}, y ajuste de pendientes log--log
        para validar la complejidad~\cite{mcgeoch2012}.
  \item Una evaluación sobre datos sintéticos controlados y sobre los corpus de referencia
        Canterbury~\cite{arnold1997} y Silesia~\cite{deorowicz2003}, con DEFLATE como punto de
        comparación externo.
\end{itemize}

\subsection{Organización del documento}
La \cref{sec:estado-arte} revisa críticamente el estado del arte: los fundamentos teóricos del
problema, las familias de algoritmos que lo resuelven y los estudios comparativos previos. La
\cref{sec:marco-teorico} presenta el marco teórico: definiciones, ambos algoritmos con su
pseudocódigo, ejemplos resueltos, el análisis de complejidad y las hipótesis H1--H4. La
\cref{sec:metodologia} describe gráficamente la metodología, la implementación, los conjuntos
de datos y el protocolo de medición. La \cref{sec:resultados} reporta los resultados de los
experimentos E1--E4, y la \cref{sec:discusion} los contrasta con la teoría. La
\cref{sec:conclusiones} responde las preguntas de investigación y plantea el trabajo futuro.
Los anexos detallan cómo reproducir el estudio y el detalle de la implementación de cada componente.

% ---------------------------------------------------------------------
%  2. Estado del arte
% ---------------------------------------------------------------------
\section{Estado del arte}
\label{sec:estado-arte}

\subsection{El problema: teoría de la información y codificación sin pérdida}
\label{sec:ea-problema}
El punto de partida de toda la disciplina es el trabajo de Shannon~\cite{shannon1948}, que
modela una fuente como un proceso aleatorio y define su entropía como el número mínimo
promedio de bits por símbolo con que puede representarse sin pérdida. El teorema de
codificación de fuente convierte la compresión en un problema bien planteado: existe un límite
inferior y la calidad de un compresor se mide por su distancia a él. Sin embargo, la entropía
depende del \emph{modelo} de la fuente: un modelo de orden cero, que trata cada byte como
independiente, ignora la redundancia contenida en secuencias repetidas, mientras que modelos
con contexto capturan dependencias entre símbolos consecutivos~\cite{bell1989}.

Bell, Witten y Cleary~\cite{bell1989} propusieron una distinción que sigue organizando el campo:
todo compresor combina un \emph{modelo} que asigna probabilidades o identifica repeticiones y
un \emph{codificador} que traduce esa información en bits. Bajo esa óptica, los métodos se
agrupan en \emph{estadísticos}, que codifican símbolos según su probabilidad, y \emph{de
diccionario}, que sustituyen secuencias por referencias. Los textos de referencia de
Sayood~\cite{sayood2017} y de Salomon y Motta~\cite{salomon2010} adoptan esta misma taxonomía, y
la revisión de Jayasankar et al.~\cite{jayasankar2021} la amplía al clasificar las técnicas por
calidad de reconstrucción, esquema de codificación, tipo de dato y aplicación. Una conclusión
común a estas revisiones es que no existe un compresor universalmente mejor: el desempeño
depende de la coincidencia entre la redundancia que explota el método y la que contienen los
datos. Esta observación motiva comparar Huffman y LZW sobre fuentes de naturaleza
deliberadamente distinta.

\subsection{Algoritmos que resuelven el problema}
\label{sec:ea-algoritmos}
\paragraph{Familia de Huffman.} Huffman~\cite{huffman1952} demostró que un algoritmo voraz que
fusiona repetidamente los dos símbolos menos probables produce un código prefijo de longitud
media mínima para una distribución conocida. Gallager~\cite{gallager1978} caracterizó estos
códigos mediante la propiedad de hermanos (\emph{sibling property}), acotó su redundancia en
función de la probabilidad del símbolo más frecuente y planteó una versión adaptativa. Esa
línea se completó con los algoritmos dinámicos de Knuth~\cite{knuth1985} y de
Vitter~\cite{vitter1987}, que actualizan el árbol tras cada símbolo y evitan la segunda pasada
y la transmisión de la tabla, a cambio de un mayor costo por símbolo. Desde el punto de vista
de la implementación, Moffat y Turpin~\cite{moffat1997} mostraron que basta transmitir las
longitudes de los códigos si se usan \emph{códigos canónicos}, que además permiten decodificar
con tablas pequeñas; la revisión de Moffat~\cite{moffat2019} sintetiza siete décadas de estas
mejoras. La limitación estructural de toda la familia es que asigna un número entero de bits a
cada símbolo, por lo que pierde hasta un bit por símbolo frente a la entropía cuando la
distribución está muy sesgada.

\paragraph{Codificación aritmética: alternativa fuera de alcance.} La codificación
aritmética~\cite{witten1987} supera esa limitación al representar el mensaje completo como un
intervalo y acercarse a la entropía del modelo con una fracción de bit de sobrecosto por
mensaje. Constituye el estándar en compresores estadísticos avanzados, pero su comparación
excede el alcance de este trabajo, centrado en los dos algoritmos clásicos de cada familia.

\paragraph{Familia Lempel--Ziv.} Ziv y Lempel propusieron dos esquemas universales que no
requieren conocer la distribución de la fuente: LZ77~\cite{ziv1977}, que codifica cada frase
como una referencia a una coincidencia previa dentro de una ventana deslizante, y
LZ78~\cite{ziv1978}, que construye incrementalmente un diccionario explícito de frases y cuya
tasa converge asintóticamente a la entropía para fuentes estacionarias y ergódicas. Storer y
Szymanski~\cite{storer1982} formalizaron la compresión por sustitución textual y propusieron
evitar referencias que no ahorran espacio, idea conocida como LZSS. Welch~\cite{welch1984}
simplificó LZ78 inicializando el diccionario con todos los símbolos del alfabeto, de modo que
la salida consiste solo en índices; el resultado, LZW, es rápido y sencillo de implementar
en hardware y software, pero su diccionario crece sin límite si no se gestiona.

\paragraph{DEFLATE: la combinación de ambas ideas.} El formato DEFLATE~\cite{rfc1951} combina
una etapa LZ77 con una codificación de Huffman canónica de los literales, longitudes y
distancias resultantes. Su éxito práctico ilustra que ambas familias son complementarias: la
etapa de diccionario elimina la repetición de frases y la etapa estadística, el sesgo en la
distribución de los símbolos resultantes. DEFLATE queda fuera del alcance de este trabajo, que
compara los dos mecanismos por separado.

\paragraph{Gestión del diccionario en LZW.} Cuando el diccionario de LZW se llena, el
compresor debe decidir qué hacer. Cheng et al.~\cite{cheng2021} comparan las políticas
habituales: el reinicio completo (\emph{reset}) es más rápido pero empeora la tasa de
compresión, mientras que el reemplazo LRU obtiene mejor tasa a costa de más tiempo y memoria;
proponen IDA-LZW, que gestiona un área de entradas poco significativas. La implementación de
este trabajo adopta el reinicio, la política más simple y predecible, y su efecto se estudia
variando el tamaño máximo del diccionario.

\subsection{Comparaciones previas y metodologías de evaluación}
\label{sec:ea-comparaciones}
La \cref{tab:trabajos-relacionados} resume los estudios que comparan directamente Huffman y
LZW o variantes de ambos. Jambek y Khairi~\cite{jambek2014} los evalúan para nodos de redes de
sensores inalámbricos y reportan que Huffman redujo los datos en alrededor de 43\,\% en
promedio y fue aproximadamente cuatro veces más rápido que LZW. Fauzan et
al.~\cite{fauzan2022}, en cambio, trabajan con archivos de texto (\texttt{.txt}, \texttt{.csv},
entre otros) y encuentran que LZW obtuvo mayor ahorro de espacio en \texttt{.txt} y
\texttt{.csv} (63.85\,\% y 77.56\,\% en promedio), con tiempos de compresión que crecen con
el tamaño del archivo. Btoush et al.~\cite{btoush2008} comparan LZW, Huffman y codificación de
longitud fija sobre textos de longitud variable, midiendo tasa, tiempo y entropía, y observan
que LZW es el peor en archivos pequeños. Bjärås~\cite{bjaras2019} evalúa cinco algoritmos
ligeros para series temporales en dispositivos IoT, incluidas variantes basadas en Huffman y en
LZW, y concluye que las basadas en Huffman ofrecen el mejor equilibrio entre tasa, memoria y
uso de CPU. Cheng et al.~\cite{cheng2021} no comparan contra Huffman, pero muestran que la
política de diccionario lleno cambia por sí sola la tasa y el tiempo de LZW.

Estos resultados, aparentemente contradictorios, son coherentes con la observación de la
\cref{sec:ea-problema}: las lecturas de sensores y las series temporales tienen poca
repetición de frases y una distribución de valores sesgada, lo que favorece a Huffman,
mientras que el texto contiene palabras y frases recurrentes, que LZW aprovecha una vez que su
diccionario se ha poblado; en archivos pequeños el diccionario no alcanza a aprender, lo que
explica la desventaja reportada por Btoush et al. Ninguno de estos trabajos, sin embargo,
aísla experimentalmente cada tipo de redundancia ni verifica cómo escalan los tiempos.

Respecto a la metodología, la literatura de algoritmia experimental ofrece pautas que los
estudios anteriores no siguen. Johnson~\cite{johnson2002} recomienda instancias de prueba
diversas y de tamaño creciente, informes reproducibles y comparaciones con la teoría.
McGeoch~\cite{mcgeoch2012} describe cómo estimar la tasa de crecimiento de un algoritmo a
partir de mediciones, por ejemplo mediante la pendiente en escala log--log. Georges et
al.~\cite{georges2007} demuestran que, en Java, la compilación JIT, la recolección de basura y
la variabilidad entre ejecuciones pueden invertir las conclusiones si se reporta una única
medición, y proponen repeticiones con intervalos de confianza. Por último, el uso de corpus
públicos como Canterbury~\cite{arnold1997} y Silesia~\cite{deorowicz2003} hace comparables los
resultados entre estudios.

\begin{table}[H]
  \centering
  \caption{Estudios comparativos relacionados.}
  \label{tab:trabajos-relacionados}
  \footnotesize
  \setlength{\tabcolsep}{4pt}
  \renewcommand{\arraystretch}{1.15}
  \begin{tabularx}{\linewidth}{@{}>{\raggedright}p{1.6cm} >{\raggedright}p{2.4cm} >{\raggedright}p{2.6cm}
                  >{\raggedright}p{2.2cm} >{\raggedright\arraybackslash}X@{}}
    \toprule
    Trabajo & Algoritmos & Datos & Métricas & Hallazgo principal \\
    \midrule
    \cite{jambek2014} & Huffman, LZW & Datos de nodos de sensores inalámbricos &
      Reducción de datos, tiempo &
      Huffman redujo los datos alrededor de 43\,\% en promedio y fue unas cuatro veces más
      rápido que LZW. \\
    \cite{fauzan2022} & Huffman, LZW & Archivos de texto (\texttt{.txt}, \texttt{.csv}, otros) &
      Ahorro de espacio, tiempo &
      LZW logró mayor ahorro en \texttt{.txt} y \texttt{.csv} (63.85\,\% y 77.56\,\% en
      promedio); el tiempo creció con el tamaño del archivo. \\
    \cite{btoush2008} & LZW, Huffman, longitud fija & Textos de longitud variable &
      Tasa, tiempo, entropía & LZW fue el peor en archivos pequeños. \\
    \cite{bjaras2019} & Cinco algoritmos ligeros (variantes Huffman y LZW, entre otros) &
      Series temporales IoT & Tasa, memoria, CPU &
      Las variantes basadas en Huffman resultaron las más equilibradas. \\
    \cite{cheng2021} & Políticas de diccionario de LZW; IDA-LZW & Conjuntos de prueba propios del estudio &
      Tasa, tiempo, memoria &
      El reinicio es más rápido pero comprime peor; LRU comprime mejor con más tiempo y
      memoria. \\
    \midrule
    Este trabajo & Huffman canónico, LZW 9--16 bits &
      Sintéticos controlados, texto real, Canterbury y Silesia &
      Tiempo (mediana, IC 95\,\%), pendiente log--log, bpb &
      Ver \cref{sec:resultados,sec:discusion}. \\
    \bottomrule
  \end{tabularx}
\end{table}

\subsection{Síntesis y brecha}
\label{sec:ea-brecha}
La revisión permite identificar tres carencias en las comparaciones previas entre Huffman y
LZW. Primero, los datos son pequeños o de un único dominio (sensores, series temporales o
texto), lo que impide separar el efecto del tipo de redundancia del efecto del algoritmo.
Segundo, ningún estudio valida la complejidad teórica: los tiempos se reportan para unos pocos
tamaños, sin estimar la tasa de crecimiento. Tercero, falta rigor estadístico: no se describen
calentamiento, repeticiones ni intervalos de confianza, pese a que en entornos con JIT esto
puede alterar las conclusiones~\cite{georges2007}.

Este trabajo aborda esas carencias con (i)~fuentes sintéticas que aíslan cada tipo de
redundancia, junto con texto real y los corpus estándar~\cite{arnold1997,deorowicz2003};
(ii)~tamaños de entrada que crecen geométricamente, de $2^{14}$ a $2^{24}$ bytes, y un
ajuste log--log contrastado con la complejidad derivada~\cite{mcgeoch2012}; y (iii)~un
protocolo con calentamiento, repeticiones, medianas e intervalos de confianza de Student,
siguiendo las recomendaciones de la algoritmia experimental~\cite{johnson2002,georges2007}.

% ---------------------------------------------------------------------
%  3. Marco teórico — D (definiciones), A (Huffman), B (LZW),
%     C (comparación teórica e hipótesis)
% ---------------------------------------------------------------------
\section{Marco teórico}
\label{sec:marco-teorico}

\subsection{Definiciones preliminares}
\label{sec:definiciones}
\begin{aporte}{D}
Sea $x = x_1x_2\cdots x_n$ una secuencia de $n$ bytes sobre el alfabeto
$\Sigma = \{0,\dots,255\}$, de la cual aparecen $\sigma \le 256$ símbolos distintos. La
frecuencia del símbolo $s$ es $f_s$ y su probabilidad empírica es $p_s = f_s/n$.

\paragraph{Fuente y entropía de orden cero.} Un modelo de orden cero trata los símbolos como
independientes e idénticamente distribuidos (i.i.d.). Su entropía empírica
\begin{equation}
  H_0(x) = -\sum_{s\,:\,f_s>0} p_s \log_2 p_s \qquad \text{[bits por símbolo]}
  \label{eq:entropia}
\end{equation}
es el número mínimo promedio de bits por símbolo de cualquier código que asigne a cada
símbolo una palabra fija~\cite{shannon1948}. Cumple $0 \le H_0 \le \log_2\sigma \le 8$, con
igualdad superior para la distribución uniforme.

\paragraph{Códigos prefijo y desigualdad de Kraft--McMillan.} Un código asigna a cada símbolo
$s$ una palabra binaria de longitud $\ell_s$; es \emph{prefijo} si ninguna palabra es prefijo
de otra, lo que permite decodificar sin separadores. Todo código unívocamente decodificable
satisface
\begin{equation}
  \sum_{s} 2^{-\ell_s} \le 1,
  \label{eq:kraft}
\end{equation}
y, recíprocamente, para cualquier conjunto de longitudes que cumpla \eqref{eq:kraft} existe un
código prefijo con esas longitudes~\cite{sayood2017,salomon2010}. La longitud media del código
es $\Lbar = \sum_s p_s \ell_s$.

\paragraph{Métricas de compresión.} Si $|C(x)|$ es el tamaño en bytes de la salida comprimida,
incluida la cabecera, se definen
\begin{equation}
  \bpb = \frac{8\,|C(x)|}{n}, \qquad
  \text{razón} = \frac{n}{|C(x)|}, \qquad
  \text{ahorro} = 1 - \frac{|C(x)|}{n}.
  \label{eq:metricas}
\end{equation}
Los bits por byte (bpb) permiten comparar directamente con $H_0$: un valor $\bpb > 8$
indica expansión.
\end{aporte}

% =====================================================================
\subsection{Codificación de Huffman}
\label{sec:huffman}
\begin{aporte}{A}
\paragraph{Idea.} Huffman~\cite{huffman1952} construye un árbol binario de abajo hacia arriba:
parte de una hoja por símbolo, ponderada por su frecuencia, y fusiona repetidamente los dos
árboles de menor peso. La profundidad de cada hoja es la longitud de su código. Se trata de un
algoritmo voraz cuya elección local ---fusionar los dos símbolos menos frecuentes--- es
óptima~\cite{cormen2022}. El \cref{alg:huffman} muestra el cálculo de longitudes tal como lo
realiza \texttt{HuffmanCodec.codeLengths}, con un montículo binario de mínimos y desempate
por orden de creación del nodo para que el resultado sea determinista.

\begin{algorithm}[H]
  \caption{Longitudes de código de Huffman (\texttt{codeLengths}).}
  \label{alg:huffman}
  \small
  \begin{algorithmic}[1]
    \Require frecuencias $f_0,\dots,f_{255}$
    \Ensure longitudes $\ell_0,\dots,\ell_{255}$ ($\ell_s = 0$ si $f_s = 0$)
    \State $Q \gets$ montículo de mínimos con una hoja por cada $s$ con $f_s > 0$, clave $w(s)=f_s$
    \If{$|Q| = 1$} \State \Return $\ell_s = 1$ para el único símbolo presente \EndIf
    \While{$|Q| > 1$}
      \State $a \gets \textsc{ExtraerMín}(Q)$;\; $b \gets \textsc{ExtraerMín}(Q)$
      \State crear nodo $z$ con $w(z) = w(a) + w(b)$;\; $\mathit{padre}(a) \gets z$;\; $\mathit{padre}(b) \gets z$
      \State $\textsc{Insertar}(Q, z)$
    \EndWhile
    \ForAll{nodo $v$ en orden decreciente de creación}
      \State $\mathit{prof}(v) \gets \mathit{prof}(\mathit{padre}(v)) + 1$ \Comment{la raíz tiene profundidad 0}
    \EndFor
    \State \Return $\ell_s = \mathit{prof}(s)$ para cada hoja $s$
  \end{algorithmic}
\end{algorithm}

\paragraph{Optimalidad y cota de entropía.} Entre todos los códigos prefijo para una
distribución dada, el de Huffman minimiza $\Lbar$~\cite{huffman1952,cormen2022}. Como su árbol
es completo (todo nodo interno tiene dos hijos), para $\sigma \ge 2$ la desigualdad
\eqref{eq:kraft} se cumple con igualdad. Además,
\begin{equation}
  H_0 \;\le\; \Lbar \;<\; H_0 + 1,
  \label{eq:cota-huffman}
\end{equation}
la cota inferior por el teorema de Shannon~\cite{shannon1948} y la superior porque el código
con longitudes $\lceil -\log_2 p_s\rceil$ satisface \eqref{eq:kraft} y Huffman no puede ser
peor~\cite{sayood2017}. Gallager~\cite{gallager1978} refinó la cota superior en función de la
probabilidad del símbolo más frecuente.

\paragraph{Códigos canónicos.} Las longitudes bastan para reconstruir el código si las palabras
se asignan de forma \emph{canónica}: en orden creciente de longitud y, dentro de cada longitud,
en orden creciente de símbolo, como especifica DEFLATE~\cite{rfc1951}. Con $c_\ell$ el número
de códigos de longitud $\ell$, el primer código de longitud $\ell$ es
\begin{equation}
  \mathit{first}_1 = 0, \qquad \mathit{first}_\ell = \left(\mathit{first}_{\ell-1} + c_{\ell-1}\right) \cdot 2,
  \label{eq:canonico}
\end{equation}
y los símbolos de esa longitud reciben códigos consecutivos a partir de él. Por eso la
implementación guarda en la cabecera solo 256 longitudes de un byte, en lugar del árbol.

\paragraph{Ejemplo: \texttt{abracadabra}.} La cadena tiene $n = 11$ y frecuencias
$a{:}5$, $b{:}2$, $r{:}2$, $c{:}1$, $d{:}1$. El algoritmo fusiona primero $c$ y $d$ (peso 2),
luego $b$ y $r$ (peso 4), después ambos subárboles (peso 6) y finalmente el resultado con $a$
(peso 11). La \cref{fig:arbol-huffman} muestra el árbol y la \cref{tab:abracadabra} las
longitudes y los códigos. El código canónico depende solo de las longitudes, por lo que puede
diferir del que se lee en el árbol sin cambiar $\Lbar$.

\begin{figure}[H]
  \centering
  \begin{tikzpicture}[
      level distance=1.25cm,
      level 1/.style={sibling distance=4.6cm},
      level 2/.style={sibling distance=2.8cm},
      level 3/.style={sibling distance=1.4cm},
      interno/.style={circle, draw=black!70, fill=gray!10, minimum size=8mm, inner sep=0pt, font=\small},
      hoja/.style={rectangle, rounded corners=2pt, draw=black!70, fill=white, minimum width=11mm,
                   minimum height=7mm, font=\small},
      edge from parent/.style={draw=black!55, -{Stealth[length=2mm]}},
      etiqueta/.style={font=\scriptsize, fill=white, inner sep=1pt}]
    \node[interno] {11}
      child { node[hoja] {\texttt{a}\,:\,5} edge from parent node[etiqueta, above left] {0} }
      child { node[interno] {6}
        child { node[interno] {2}
          child { node[hoja] {\texttt{c}\,:\,1} edge from parent node[etiqueta, above left] {0} }
          child { node[hoja] {\texttt{d}\,:\,1} edge from parent node[etiqueta, above right] {1} }
          edge from parent node[etiqueta, above left] {0} }
        child { node[interno] {4}
          child { node[hoja] {\texttt{b}\,:\,2} edge from parent node[etiqueta, above left] {0} }
          child { node[hoja] {\texttt{r}\,:\,2} edge from parent node[etiqueta, above right] {1} }
          edge from parent node[etiqueta, above right] {1} }
        edge from parent node[etiqueta, above right] {1} };
  \end{tikzpicture}
  \caption{Árbol de Huffman para \texttt{abracadabra}; los nodos internos muestran el peso acumulado.}
  \label{fig:arbol-huffman}
\end{figure}

\begin{table}[H]
  \centering
  \caption{Longitudes y códigos para \texttt{abracadabra}.}
  \label{tab:abracadabra}
  \small
  \begin{tabular}{@{}c c c c c@{}}
    \toprule
    Símbolo & Frecuencia & Longitud $\ell_s$ & Código del árbol & Código canónico \\
    \midrule
    \texttt{a} & 5 & 1 & \texttt{0}   & \texttt{0}   \\
    \texttt{b} & 2 & 3 & \texttt{110} & \texttt{100} \\
    \texttt{c} & 1 & 3 & \texttt{100} & \texttt{101} \\
    \texttt{d} & 1 & 3 & \texttt{101} & \texttt{110} \\
    \texttt{r} & 2 & 3 & \texttt{111} & \texttt{111} \\
    \bottomrule
  \end{tabular}
\end{table}

Con $c_1 = 1$, $c_2 = 0$ y $c_3 = 4$, la \cref{eq:canonico} da $\mathit{first}_1 = 0$,
$\mathit{first}_2 = (0 + 1)\cdot 2 = 2$ y $\mathit{first}_3 = (2 + 0)\cdot 2 = 4 = \texttt{100}_2$, de modo que $b, c, d, r$
reciben \texttt{100}, \texttt{101}, \texttt{110} y \texttt{111}. La cadena se codifica como
\texttt{0\,100\,111\,0\,101\,0\,110\,0\,100\,111\,0}, es decir, 23 bits frente a 88 bits sin
comprimir. Así, $\Lbar = 23/11 \approx 2.091$ y $H_0 \approx 2.040$, lo que verifica
\eqref{eq:cota-huffman}; además $2^{-1} + 4\cdot 2^{-3} = 1$, la igualdad de Kraft.

\paragraph{Decodificación canónica.} Siguiendo a Moffat y Turpin~\cite{moffat1997}, el
decodificador no reconstruye el árbol: lee un bit por iteración y, para cada longitud $\ell$,
comprueba si el valor acumulado cae en el rango de los $c_\ell$ códigos de esa longitud
(\cref{alg:huffman-dec}). Los símbolos se almacenan ordenados por código canónico, de modo que
el índice del símbolo se obtiene por simple resta.

\begin{algorithm}[H]
  \caption{Decodificación canónica de un símbolo (\texttt{decodeSymbol}).}
  \label{alg:huffman-dec}
  \small
  \begin{algorithmic}[1]
    \Require lector de bits; $c_\ell$ por longitud; símbolos ordenados por código canónico
    \State $\mathit{code} \gets 0$;\; $\mathit{first} \gets 0$;\; $\mathit{index} \gets 0$
    \For{$\ell = 1$ \textbf{hasta} $L_{\max}$}
      \State $\mathit{code} \gets \mathit{code} \mathbin{|} \textsc{LeerBit}()$
      \If{$\mathit{code} - \mathit{first} < c_\ell$}
        \State \Return $\textit{símbolos}[\mathit{index} + \mathit{code} - \mathit{first}]$
      \EndIf
      \State $\mathit{index} \gets \mathit{index} + c_\ell$;\;
             $\mathit{first} \gets (\mathit{first} + c_\ell)\cdot 2$;\;
             $\mathit{code} \gets \mathit{code}\cdot 2$
    \EndFor
    \State \textbf{error}: código inválido
  \end{algorithmic}
\end{algorithm}

\paragraph{Análisis de complejidad.} Sea $n$ el tamaño de la entrada y $\sigma \le 256$ el
número de símbolos presentes.
\begin{itemize}
  \item \emph{Conteo de frecuencias:} una pasada con un incremento por byte, $\Theta(n)$.
  \item \emph{Construcción del código:} $\sigma$ inserciones iniciales y $\sigma - 1$
        iteraciones con dos extracciones y una inserción en el montículo, cada una
        $O(\log\sigma)$; en total $O(\sigma\log\sigma)$. La asignación de profundidades y de
        códigos canónicos cuesta $O(\sigma + L_{\max})$.
  \item \emph{Codificación:} segunda pasada sobre la entrada; cada símbolo escribe $\ell_s$
        bits con \texttt{BitWriter} en tiempo $O(1 + \ell_s/8)$, de modo que el total es
        $\Theta(n)$. El costo global de codificar es $\Theta(n + \sigma\log\sigma)$, es decir,
        $\Theta(n)$ para un alfabeto fijo.
  \item \emph{Decodificación:} el \cref{alg:huffman-dec} realiza $\ell_s$ iteraciones por
        símbolo, por lo que el total es $\sum_i \ell_{x_i} = n\Lbar$ iteraciones, es decir,
        $\Theta(n\Lbar)$. Por \eqref{eq:cota-huffman}, $\Lbar < H_0 + 1 \le \log_2 256 + 1 = 9$,
        así que el costo es lineal en $n$, pero con una constante proporcional a la longitud
        media del código y, por tanto, a la entropía de la fuente.
  \item \emph{Espacio adicional:} $O(\sigma)$ para el montículo, los pesos, los padres y las
        tablas canónicas, sin contar los arreglos de entrada y salida. La cabecera ocupa
        $2 + 4 + 256 = 262$ bytes, un sobrecosto de $8\cdot 262/n$ bpb que solo es apreciable
        en entradas pequeñas.
\end{itemize}
\end{aporte}

% =====================================================================
\subsection{Algoritmo Lempel--Ziv--Welch (LZW)}
\label{sec:lzw}
\begin{aporte}{B}
\paragraph{Idea.} LZW~\cite{welch1984} es una variante de LZ78~\cite{ziv1978} que inicializa el
diccionario con los 256 bytes posibles. El codificador extiende la frase actual $w$ mientras
$w\,c$ siga en el diccionario; cuando no lo está, emite el código de $w$, agrega $w\,c$ como
nueva entrada y reinicia la frase con $c$. El decodificador reconstruye el mismo diccionario
a partir de los códigos, sin que este se transmita. En esta implementación el código 256 se
reserva para \texttt{CLEAR}, de modo que las frases nuevas comienzan en 257.

\paragraph{Códigos de ancho variable y reinicio.} Los códigos se escriben con el menor ancho
capaz de representar todos los códigos posibles en ese momento,
$\mathit{ancho}(s) = \max\left(9, \lceil\log_2 s\rceil\right)$, donde $s$ es el siguiente
código libre; el ancho crece de 9 hasta $b$ bits, con $b \in [9, 16]$ configurable
(\texttt{maxBits}, 16 por defecto). Cuando el diccionario alcanza $2^b$ entradas, el
codificador emite \texttt{CLEAR} y ambos lados vuelven al diccionario inicial. Esta política
de reinicio mantiene el modelo adaptativo ante cambios en los datos y acota la memoria, a costa
de olvidar frases útiles~\cite{cheng2021}. Los \cref{alg:lzw-enc,alg:lzw-dec} resumen
\texttt{LzwCodec.compress} y \texttt{LzwCodec.decode}.

\begin{algorithm}[H]
  \caption{Codificación LZW con reinicio (\texttt{compress}).}
  \label{alg:lzw-enc}
  \small
  \begin{algorithmic}[1]
    \Require bytes $x_1\cdots x_n$ ($n \ge 1$), ancho máximo $b$
    \State $D \gets$ diccionario con los 256 bytes;\; $s \gets 257$;\; $w \gets x_1$
    \For{$i = 2$ \textbf{hasta} $n$}
      \State $c \gets x_i$
      \If{$(w, c) \in D$}
        \State $w \gets D[w, c]$ \Comment{extender la frase}
      \Else
        \State emitir $w$ con $\mathit{ancho}(s)$ bits
        \If{$s < 2^b$}
          \State $D[w, c] \gets s$;\; $s \gets s + 1$
        \Else
          \State emitir \texttt{CLEAR} con $\mathit{ancho}(s)$ bits;\; vaciar $D$;\; $s \gets 257$
        \EndIf
        \State $w \gets c$
      \EndIf
    \EndFor
    \State emitir $w$ con $\mathit{ancho}(s)$ bits
  \end{algorithmic}
\end{algorithm}

\begin{algorithm}[H]
  \caption{Decodificación LZW (\texttt{decode}).}
  \label{alg:lzw-dec}
  \small
  \begin{algorithmic}[1]
    \Require flujo de códigos, longitud original $n$, ancho máximo $b$
    \State $s \gets 257$;\; $\mathit{prev} \gets \bot$
    \While{se hayan producido menos de $n$ bytes}
      \State leer $k$ con $\mathit{ancho}(s)$ bits si $\mathit{prev} = \bot$, o con
             $\mathit{ancho}(\min(s + 1, 2^b))$ en otro caso
      \If{$k = \texttt{CLEAR}$} \State $s \gets 257$;\; $\mathit{prev} \gets \bot$;\; \textbf{continuar} \EndIf
      \If{$\mathit{prev} = \bot$} \State producir el byte $k$;\; $\mathit{prev} \gets k$;\; \textbf{continuar} \EndIf
      \If{$s < 2^b$}
        \State $h \gets \mathit{primero}(k)$ si $k < s$; si no, $h \gets \mathit{primero}(\mathit{prev})$ \Comment{caso $KwKwK$}
        \State $D[s] \gets (\mathit{prev}, h)$;\; $s \gets s + 1$
      \EndIf
      \State producir la frase $D[k]$ recorriendo prefijos de atrás hacia adelante;\; $\mathit{prev} \gets k$
    \EndWhile
  \end{algorithmic}
\end{algorithm}

El decodificador agrega cada entrada un código más tarde que el codificador, porque necesita
el primer byte de la frase siguiente; por eso lee con el ancho correspondiente a $s + 1$.

\paragraph{Ejemplo: \texttt{TOBEORNOTTOBEORTOBEORNOT}.} La \cref{tab:lzw-welch} reproduce el
ejemplo clásico de Welch con la numeración de esta implementación. Las 24 letras se codifican
con 16 códigos, la secuencia
\texttt{T O B E O R N O T 257 259 261 266 260 262 264}, que es exactamente la que verifica la
prueba \texttt{emitsTextbookCodeSequence}. Como el diccionario nunca supera 512 entradas, todos
los códigos usan 9 bits: $16 \cdot 9 = 144$ bits frente a $24 \cdot 8 = 192$ bits de entrada.

\begin{table}[H]
  \centering
  \caption{Codificación LZW de \texttt{TOBEORNOTTOBEORTOBEORNOT} (solo los pasos que emiten).}
  \label{tab:lzw-welch}
  \small
  \begin{tabular}{@{}r l c r l@{}}
    \toprule
    Paso & Frase $w$ & Siguiente $c$ & Código emitido & Nueva entrada \\
    \midrule
     1 & \texttt{T}   & \texttt{O} & 84 (\texttt{T}) & 257 = \texttt{TO} \\
     2 & \texttt{O}   & \texttt{B} & 79 (\texttt{O}) & 258 = \texttt{OB} \\
     3 & \texttt{B}   & \texttt{E} & 66 (\texttt{B}) & 259 = \texttt{BE} \\
     4 & \texttt{E}   & \texttt{O} & 69 (\texttt{E}) & 260 = \texttt{EO} \\
     5 & \texttt{O}   & \texttt{R} & 79 (\texttt{O}) & 261 = \texttt{OR} \\
     6 & \texttt{R}   & \texttt{N} & 82 (\texttt{R}) & 262 = \texttt{RN} \\
     7 & \texttt{N}   & \texttt{O} & 78 (\texttt{N}) & 263 = \texttt{NO} \\
     8 & \texttt{O}   & \texttt{T} & 79 (\texttt{O}) & 264 = \texttt{OT} \\
     9 & \texttt{T}   & \texttt{T} & 84 (\texttt{T}) & 265 = \texttt{TT} \\
    10 & \texttt{TO}  & \texttt{B} & 257 & 266 = \texttt{TOB} \\
    11 & \texttt{BE}  & \texttt{O} & 259 & 267 = \texttt{BEO} \\
    12 & \texttt{OR}  & \texttt{T} & 261 & 268 = \texttt{ORT} \\
    13 & \texttt{TOB} & \texttt{E} & 266 & 269 = \texttt{TOBE} \\
    14 & \texttt{EO}  & \texttt{R} & 260 & 270 = \texttt{EOR} \\
    15 & \texttt{RN}  & \texttt{O} & 262 & 271 = \texttt{RNO} \\
    16 & \texttt{OT}  & (fin)      & 264 & --- \\
    \bottomrule
  \end{tabular}
\end{table}

\paragraph{El caso especial $KwKwK$.} Si la entrada contiene $K\,w\,K\,w\,K$, con $K$ un byte y
$w$ una cadena, el codificador puede emitir un código que acaba de crear y que el decodificador
aún no conoce. Con la entrada \texttt{aaa}, por ejemplo, el codificador emite 97
(\texttt{a}), agrega $257 = \texttt{aa}$ y luego emite 257. Al recibir 257, el decodificador
tiene $s = 257$: el código es igual al siguiente libre. La única frase consistente es la
anterior seguida de su propio primer byte, $\texttt{a} + \texttt{a} = \texttt{aa}$, que es lo
que hace la rama $k = s$ del \cref{alg:lzw-dec}; la prueba \texttt{handlesCodeNotYetInDictionary}
cubre este caso.

\paragraph{Análisis de complejidad.} Sea $D = 2^b$ el tamaño máximo del diccionario.
\begin{itemize}
  \item \emph{Codificación:} el diccionario se implementa como una tabla hash de
        direccionamiento abierto con sondeo lineal y hashing de Fibonacci (multiplicación por
        \texttt{0x9E3779B1}), indexada por la clave $(w, c)$. Su capacidad es $2D$, de modo que
        el factor de carga nunca supera $\alpha = 0.5$. Cada byte de entrada provoca una
        búsqueda; bajo la hipótesis de hashing uniforme, una búsqueda fallida requiere a lo
        sumo $1/(1-\alpha) = 2$ sondeos esperados~\cite{cormen2022}. Cada código emitido cuesta
        $O(1)$ adicional. Un reinicio limpia la tabla en $O(D)$, pero ocurre como mucho una vez
        cada $D - 257$ códigos, y cada código consume al menos un byte, de modo que su costo
        amortizado es $O(1)$ por byte. En total, $\Theta(n)$ esperado.
  \item \emph{Decodificación:} cada código se resuelve con accesos a arreglos
        (\texttt{prefix}, \texttt{suffix}, \texttt{firstByte}, \texttt{phraseLength}) y cada
        byte de salida se escribe exactamente una vez al recorrer la frase hacia atrás; el
        costo es $\Theta(n)$ en el peor caso, sin dependencia de una función hash.
  \item \emph{Espacio adicional:} $O(D)$, independiente de $n$. Para $b = 16$, el codificador
        usa dos arreglos de $2^{17}$ enteros y el decodificador cuatro arreglos de $2^{16}$
        entradas. La cabecera ocupa solo 7 bytes (magia, $b$ y longitud original).
\end{itemize}
\end{aporte}

% =====================================================================
\subsection{Comparación teórica}
\label{sec:comparacion-teorica}
\begin{aporte}{C}
La \cref{tab:comparacion-teorica} resume las diferencias estructurales entre ambos
algoritmos tal como se implementaron. La diferencia esencial no es asintótica ---ambos son
lineales--- sino de \emph{modelo}: Huffman es estático y de dos pasadas, y explota solo la
distribución de frecuencias de los bytes; LZW es adaptativo y de una pasada, y explota
frases repetidas, pero paga al menos 9 bits por código emitido.

\begin{table}[H]
  \centering
  \caption{Comparación teórica de las implementaciones.}
  \label{tab:comparacion-teorica}
  \small
  \renewcommand{\arraystretch}{1.15}
  \begin{tabularx}{\linewidth}{@{}>{\raggedright}p{3.4cm} >{\raggedright}X >{\raggedright\arraybackslash}X@{}}
    \toprule
    Aspecto & Huffman canónico & LZW (9--$b$ bits, \texttt{CLEAR}) \\
    \midrule
    Tiempo de codificación & $\Theta(n + \sigma\log\sigma) = \Theta(n)$ & $\Theta(n)$ esperado (tabla hash) \\
    Tiempo de decodificación & $\Theta(n\Lbar)$, con $\Lbar < H_0 + 1 \le 9$ & $\Theta(n)$ en el peor caso \\
    Espacio adicional & $O(\sigma)$ & $O(2^b)$, independiente de $n$ \\
    Modelo & Estático, dos pasadas; requiere transmitir el código & Adaptativo, una pasada; el diccionario no se transmite \\
    Redundancia explotada & Frecuencias de símbolos (orden cero) & Frases repetidas \\
    Garantía de tasa & $H_0 \le \Lbar < H_0 + 1$ & Asintótica para fuentes estacionarias~\cite{ziv1978}; sin cota para $n$ finito \\
    Cabecera & 262 bytes & 7 bytes \\
    \bottomrule
  \end{tabularx}
\end{table}
\end{aporte}

\subsection{Hipótesis}
\label{sec:hipotesis}
\begin{aporte}{C}
Del análisis anterior se derivan cuatro hipótesis que guían los experimentos de la
\cref{sec:metodologia}:
\begin{enumerate}[label=\textbf{H\arabic*.}, leftmargin=*]
  \item \label{hip:h1} Los codificadores y decodificadores de ambos algoritmos escalan
        linealmente con el tamaño de la entrada: la pendiente del ajuste log--log del tiempo
        frente a $n$ es cercana a 1.
  \item \label{hip:h2} Los bits por byte de Huffman se ubican en $[H_0, H_0 + 1)$ y, en fuentes
        i.i.d., quedan cerca de $H_0$.
  \item \label{hip:h3} LZW supera a Huffman en datos con frases repetidas (texto natural y
        sintético), pierde frente a él en datos i.i.d. y expande los datos aleatorios
        uniformes ($\bpb > 8$).
  \item \label{hip:h4} El tiempo de decodificación de Huffman crece con la longitud media del
        código (aproximadamente, con la entropía), mientras que el de LZW es prácticamente
        insensible a ella.
\end{enumerate}
\end{aporte}

% ---------------------------------------------------------------------
%  4. Metodología — C (general), A y B (implementación y pruebas propias)
% ---------------------------------------------------------------------
\section{Metodología}
\label{sec:metodologia}

\subsection{Proceso de comparación}
\label{sec:proceso}
La \cref{fig:metodologia} resume los pasos de la comparación. Se parte de conjuntos de datos
elegidos para aislar cada tipo de redundancia; ambos algoritmos se implementan detrás de una
misma interfaz y se validan antes de medir; cada configuración se mide tras calentar el
compilador JIT, con repeticiones cronometradas y resúmenes estadísticos robustos; finalmente,
la tasa de crecimiento estimada y la tasa de compresión se contrastan con las hipótesis
H1--H4 derivadas del análisis teórico. El diseño sigue las pautas de la algoritmia
experimental~\cite{johnson2002,mcgeoch2012} y de la evaluación de rendimiento en
Java~\cite{georges2007}.

\begin{figure}[H]
  \centering
  \begin{tikzpicture}[
      node distance=0.55cm and 0.62cm,
      paso/.style={rectangle, rounded corners=3pt, draw=black!70, fill=gray!8,
                   text width=2.7cm, minimum height=1.5cm, align=center, font=\footnotesize},
      datos/.style={cylinder, shape border rotate=90, aspect=0.2, draw=black!70,
                    fill=gray!15, text width=2.35cm, minimum height=1.7cm, align=center,
                    font=\footnotesize},
      flecha/.style={-{Stealth[length=2.2mm]}, thick, draw=black!75},
      fase/.style={draw=gray!55, dashed, rounded corners=5pt, inner sep=5pt},
      etiq/.style={font=\scriptsize\bfseries\sffamily, text=gray!65!black}]
    \node[datos] (d) {Conjuntos de datos: sintéticos, texto real, corpus};
    \node[paso, right=of d] (i) {Implementación (\texttt{Codec}): Huffman y LZW};
    \node[paso, right=of i] (v) {Validación: 43 pruebas JUnit e ida y vuelta};
    \node[paso, right=of v] (w) {Calentamiento global y por configuración ($W$)};
    \node[paso, below=1.5cm of w] (r) {$R$ ejecuciones cronometradas: codificar y decodificar};
    \node[paso, left=of r] (s) {Estadística: mediana e IC 95\,\% (t de Student)};
    \node[paso, left=of s] (f) {Ajuste log--log: pendiente y $R^2$; bpb frente a $H_0$};
    \node[paso, left=of f, fill=gray!22] (c) {Contraste con la teoría: H1--H4};
    \draw[flecha] (d) -- (i);
    \draw[flecha] (i) -- (v);
    \draw[flecha] (v) -- (w);
    \draw[flecha] (w) -- (r);
    \draw[flecha] (r) -- (s);
    \draw[flecha] (s) -- (f);
    \draw[flecha] (f) -- (c);
    \draw[flecha, dashed] (v.north) -- ++(0,0.45) -| node[pos=0.25, above, font=\scriptsize]
      {falla: corregir} (i.north);
    % Coordenadas auxiliares: un punto calc dentro de fit= rompe el analizador de claves.
    \coordinate (vtope) at ($(v.north)+(0,0.75)$);
    \coordinate (wtope) at ($(w.north)+(0,0.75)$);
    \node[fase, fit=(d)(i)(v)(vtope),
          label={[anchor=south west, etiq]north west:PREPARACIÓN}] {};
    \node[fase, fit=(w)(r)(wtope),
          label={[anchor=south west, etiq]north west:MEDICIÓN}] {};
    \node[fase, fit=(c)(f)(s),
          label={[anchor=north west, etiq]south west:ANÁLISIS}] {};
  \end{tikzpicture}
  \caption{Metodología de la comparación entre Huffman y LZW.}
  \label{fig:metodologia}
\end{figure}

\subsection{Implementación}
\label{sec:implementacion}
\paragraph{Arquitectura.} El código es un proyecto NetBeans de tipo \emph{Java with Ant}
compilado para Java~17 y ejecutado con el JDK~21. Todos los algoritmos implementan la interfaz
\texttt{Codec} (métodos \texttt{name}, \texttt{encode} y \texttt{decode}), que exige que
\texttt{decode(encode(x))} sea igual a \texttt{x} y que los datos corruptos se rechacen con
\texttt{IllegalArgumentException}. Gracias a esa interfaz, el arnés de medición trata por igual
a Huffman y a LZW. La \cref{tab:paquetes} describe los paquetes.

\begin{table}[H]
  \centering
  \caption{Paquetes del proyecto (\texttt{compression}).}
  \label{tab:paquetes}
  \small
  \begin{tabularx}{\linewidth}{@{}l >{\raggedright}p{3.9cm} >{\raggedright\arraybackslash}X@{}}
    \toprule
    Paquete & Clases & Responsabilidad \\
    \midrule
    (raíz) & \texttt{Codec} & Interfaz común de los algoritmos. \\
    \texttt{io} & \texttt{BitWriter}, \texttt{BitReader} & E/S de campos de bits de ancho variable, bit más significativo primero, con acumulador de 64 bits. \\
    \texttt{huffman} & \texttt{HuffmanCodec} & Longitudes con montículo, códigos canónicos y decodificación canónica. \\
    \texttt{lzw} & \texttt{LzwCodec} & Códigos de 9 a 16 bits, diccionario con tabla hash y reinicio con \texttt{CLEAR}. \\
    \texttt{bench} & \texttt{BenchmarkRunner}, \texttt{Stats}, \texttt{DataGenerator}, \texttt{ResultWriter} & Experimentos E1--E4, estadística, generación de datos y escritura de CSV y macros \LaTeX. \\
    \texttt{app} & \texttt{Main} & Interfaz de línea de comandos (\texttt{compress}, \texttt{decompress}, \texttt{stats}, \texttt{demo}). \\
    \bottomrule
  \end{tabularx}
\end{table}

\paragraph{Núcleo de Huffman.} El \cref{lst:huffman} muestra el bucle principal de
\ruta{HuffmanCodec.codeLengths}. Los nodos internos reciben índices crecientes a partir de
256, de modo que un padre siempre tiene un índice mayor que sus hijos y basta un recorrido
descendente para asignar profundidades sin recursión. El montículo es un
\texttt{PriorityQueue} que desempata por índice, lo que hace el código determinista.

\begin{lstlisting}[caption={Fusión de nodos y cálculo de profundidades en \texttt{HuffmanCodec.codeLengths}.},label={lst:huffman}]
int next = ALPHABET; // los nodos internos reciben los índices 256, 257, ...
while (heap.size() > 1) {
    int a = heap.poll();
    int b = heap.poll();
    weight[next] = weight[a] + weight[b];
    parent[a] = next;
    parent[b] = next;
    heap.add(next++);
}

// Un padre siempre tiene un índice mayor que sus hijos: basta un recorrido descendente.
int root = next - 1;
int[] depth = new int[maxNodes];
for (int node = root - 1; node >= ALPHABET; node--) {
    depth[node] = depth[parent[node]] + 1;
}
\end{lstlisting}

\paragraph{Núcleo de LZW.} El \cref{lst:lzw} reproduce \texttt{LzwCodec.compress}. La clave de
la tabla hash concatena el código de la frase y el byte siguiente en un entero; una búsqueda
exitosa extiende la frase, y una fallida emite el código y deja preparada la ranura
(\texttt{slot}) donde se inserta la nueva entrada, sin una segunda búsqueda. El destino de los
códigos (\texttt{CodeSink}) es el \texttt{BitWriter} al comprimir, o un arreglo al trazar los
códigos en las pruebas.

\begin{lstlisting}[caption={Bucle de codificación de \texttt{LzwCodec.compress}.},label={lst:lzw}]
private void compress(byte[] input, CodeSink sink) {
    if (input.length == 0) {
        return;
    }
    int maxSize = 1 << maxBits;
    PhraseTable table = new PhraseTable(maxBits);
    int size = FIRST_FREE;
    int phrase = input[0] & 0xFF;
    for (int i = 1; i < input.length; i++) {
        int next = input[i] & 0xFF;
        int key = (phrase << 8) | next;
        int slot = table.find(key);
        if (table.isOccupied(slot)) {
            phrase = table.valueAt(slot); // extender la frase actual
            continue;
        }
        sink.put(phrase, width(size));
        if (size < maxSize) {
            table.insert(slot, key, size++);
        } else {
            sink.put(CLEAR, width(size));
            table.clear();
            size = FIRST_FREE;
        }
        phrase = next;
    }
    sink.put(phrase, width(size));
}
\end{lstlisting}

\subsection{Validación de la corrección}
\label{sec:pruebas}
Antes de medir, la corrección se verifica en dos niveles. El primero es una batería de 43
pruebas JUnit~4 que se ejecuta con \texttt{ant test}. El segundo está incorporado en el arnés:
cada configuración medida se codifica y decodifica una vez y, si el resultado no es idéntico
a la entrada, la campaña se detiene con un error. Así, ningún tiempo reportado corresponde a
una ejecución incorrecta. La \cref{tab:pruebas} resume la batería por clase.

\begin{table}[H]
  \centering
  \caption{Batería de pruebas JUnit (43 pruebas).}
  \label{tab:pruebas}
  \small
  \begin{tabular}{@{}l c p{9.4cm}@{}}
    \toprule
    Clase de prueba & Pruebas & Aspectos verificados \\
    \midrule
    \texttt{HuffmanCodecTest} & 11 & Ida y vuelta (vacío, un byte, un solo símbolo repetido con 1 bit por símbolo, los 256 valores, datos aleatorios y sesgados, texto); longitudes óptimas del ejemplo de libro; igualdad de Kraft; códigos canónicos libres de prefijo y ordenados por longitud; rechazo de entradas corruptas; nombre del algoritmo. \\
    \texttt{BitStreamTest} & 5 & Orden MSB primero y relleno con ceros; anchos aleatorios hasta 64 bits; desplazamiento inicial; lectura más allá del final; anchos inválidos. \\
    \texttt{LzwCodecTest} & 8 & Secuencia del ejemplo de Welch; vacío y un byte; caso $KwKwK$; todos los anchos de 9 a 16 bits con reinicios del diccionario; los 256 valores; compresión de datos repetitivos; configuración inválida y datos corruptos; ancho registrado en la cabecera. \\
    \texttt{StatsTest} & 6 & Resumen de muestras, mediana con cantidad par, cuantil normal para muestras grandes, rechazo de muestras insuficientes, exponente de la ley de potencia en escala log--log, entropía de orden cero. \\
    \texttt{DataGeneratorTest} & 3 & Determinismo por semilla, entropía esperada de cada fuente, uso exacto de $k$ símbolos. \\
    \texttt{BenchmarkRunnerTest} & 4 & Una corrida rápida escribe todos los archivos de resultados; \texttt{--only e2} escribe solo el barrido; argumentos desconocidos o incompletos. \\
    \texttt{MainTest} & 6 & Compresión y descompresión de archivos con ambos algoritmos; comando \texttt{stats}; algoritmo y comando desconocidos (código de salida 2); archivo corrupto (código 1); demostración sin argumentos. \\
    \bottomrule
  \end{tabular}
\end{table}

\subsection{Conjuntos de datos}
\label{sec:datos}
Para separar el efecto del algoritmo del efecto del tipo de redundancia se combinan fuentes
sintéticas, cuyas propiedades se conocen por construcción, con texto real y con los corpus
estándar de evaluación de compresores (\cref{tab:datos}). Las fuentes sintéticas se generan
con \texttt{java.util.Random} y semilla fija 2026, por lo que son reproducibles.

\begin{table}[H]
  \centering
  \caption{Conjuntos de datos.}
  \label{tab:datos}
  \small
  \renewcommand{\arraystretch}{1.12}
  \begin{tabularx}{\linewidth}{@{}>{\raggedright}p{2.6cm} >{\raggedright}X >{\raggedright}p{4.0cm} >{\raggedright\arraybackslash}p{1.4cm}@{}}
    \toprule
    Conjunto & Construcción & Redundancia que contiene & Uso \\
    \midrule
    \texttt{uniform} & Bytes i.i.d. uniformes en $\{0,\dots,255\}$ & Ninguna ($H_0 \approx 8$) & E1 \\
    \texttt{skewed} & Símbolos i.i.d. geométricos con $p = 0.15$, truncados a 255 & Frecuencias sesgadas, sin frases & E1 \\
    \texttt{synthtext} & Palabras de un vocabulario de 2000 (2 a 9 letras con frecuencias del inglés) elegidas con ley de Zipf de exponente 1.1, separadas por espacios y saltos de línea & Frases repetidas y frecuencias sesgadas & E1 \\
    \texttt{dickens} & Prefijos del archivo \texttt{dickens} de Silesia (texto en inglés) & Texto natural & E1, E4 \\
    Alfabetos $k$ & Símbolos i.i.d. uniformes en $\{0,\dots,k-1\}$, $k = 2, 4, \dots, 256$, $n = 2^{23}$ & Entropía controlada $H_0 = \log_2 k$ & E2 \\
    Canterbury~\cite{arnold1997} & 11 archivos del corpus estándar y 3 del corpus grande (\texttt{E.coli}, \texttt{bible.txt}, \texttt{world192.txt}) & Texto, código, binarios, ADN & E3, E4 \\
    Silesia~\cite{deorowicz2003} & 12 archivos de tipos diversos (texto, XML, bases de datos, ejecutables, imágenes médicas) & Mezcla realista & E3, E4 \\
    \bottomrule
  \end{tabularx}
\end{table}

Los corpus se descargan con \texttt{data/download-corpora.sh}, que verifica sumas SHA-256 antes
de extraerlos.

\subsection{Métricas}
\label{sec:metricas}
\begin{itemize}
  \item \emph{Tiempo:} mediana, en milisegundos, de $R$ ejecuciones de codificación y de
        decodificación, medidas por separado con \texttt{System.nanoTime}.
  \item \emph{Dispersión:} semiancho del intervalo de confianza del 95\,\% de la media,
        $t_{0.975,\,R-1}\, s/\sqrt{R}$, con la tabla de cuantiles de Student para hasta 30 grados
        de libertad y $1.96$ por encima.
  \item \emph{Costo por byte:} $\text{ns/byte} = 10^6 \cdot \text{mediana [ms]}/n$; a
        complejidad lineal, esta métrica debe ser aproximadamente constante en $n$. En el corpus
        se reporta el rendimiento en MB/s, $n/10^6$ dividido entre la mediana en segundos.
  \item \emph{Compresión:} bits por byte según la \cref{eq:metricas}, incluida la cabecera, y su
        comparación con $H_0$ de la \cref{eq:entropia}.
  \item \emph{Tasa de crecimiento:} pendiente $\hat{k}$ y $R^2$ del ajuste por mínimos cuadrados
        de $\ln t = \hat{k}\ln n + c$; para un algoritmo $\Theta(n^k)$ la pendiente tiende a
        $k$~\cite{mcgeoch2012}. El ajuste usa los tamaños $n \ge 2^{16}$, para que los costos
        fijos no dominen.
\end{itemize}

\subsection{Protocolo de medición}
\label{sec:protocolo}
El protocolo implementado en \texttt{BenchmarkRunner} responde a los riesgos que Georges et
al.~\cite{georges2007} identifican en la medición de programas Java ---compilación JIT,
recolección de basura y eliminación de código muerto--- y a las recomendaciones de
reproducibilidad de Johnson~\cite{johnson2002} y McGeoch~\cite{mcgeoch2012}:
\begin{enumerate}[leftmargin=*]
  \item \emph{Calentamiento global:} antes del primer experimento, cada códec (Huffman y LZW de
        16 bits) codifica y decodifica tres veces 4~MiB de cada fuente sintética, para
        que los primeros puntos no se midan en el intérprete.
  \item \emph{Preparación de cada configuración:} se invoca \texttt{System.gc()}, se codifica una
        vez, se registra el tamaño comprimido y se verifica la ida y vuelta.
  \item \emph{Calentamiento local y repeticiones:} se ejecutan $W$ iteraciones no medidas y luego
        $R$ iteraciones medidas, según el tamaño: $W = 5$, $R = 21$ para $n \le 1$~MiB;
        $W = 3$, $R = 11$ para $n \le 16$~MiB; y $W = 1$, $R = 7$ por encima.
  \item \emph{Sumidero:} las salidas se acumulan en un campo \texttt{volatile}
        (\texttt{blackhole}) para impedir que el JIT elimine el trabajo medido.
  \item \emph{Reproducibilidad:} semilla fija, parámetros de la JVM fijos
        y escritura automática de los resultados en CSV y en macros \LaTeX{} que este informe
        importa sin transcripción manual. Los parámetros de la JVM son \texttt{-Xms3g},
        \texttt{-Xmx3g} y \texttt{-XX:+UseParallelGC}.
\end{enumerate}

\subsection{Diseño experimental}
\label{sec:experimentos}
La \cref{tab:experimentos} relaciona cada experimento con la variable que manipula, los
parámetros fijos, las salidas medidas y la hipótesis que contrasta.

\begin{table}[H]
  \centering
  \caption{Experimentos.}
  \label{tab:experimentos}
  \small
  \renewcommand{\arraystretch}{1.12}
  \begin{tabularx}{\linewidth}{@{}>{\raggedright}p{0.75cm} >{\raggedright}p{2.9cm} >{\raggedright}X >{\raggedright}p{3.7cm} >{\raggedright\arraybackslash}p{1.9cm}@{}}
    \toprule
    Exp. & Variable & Parámetros fijos & Salidas & Hipótesis \\
    \midrule
    E1 & $n = 2^{14},\dots,2^{24}$ (\texttt{dickens}: hasta $2^{23}$) &
      LZW con $b = 16$; fuentes \texttt{uniform}, \texttt{skewed}, \texttt{synthtext}, \texttt{dickens} &
      Tiempos (mediana, IC), ns/byte, bpb, pendiente y $R^2$ & H1 (y H2, H3) \\
    E2 & $k = 2, 4, \dots, 256$ ($H_0 = \log_2 k$) &
      $n = 2^{23}$; fuente i.i.d. uniforme en $k$ símbolos; LZW con $b = 16$ &
      bpb, tiempos de codificación y decodificación & H2, H4 (y H3) \\
    E3 & Archivo del corpus &
      Canterbury y Silesia completos; LZW con $b = 16$ &
      bpb, MB/s, archivos donde gana LZW & H3 (y H2) \\
    E4 & $b = 9, \dots, 16$ &
      \texttt{dickens} y \texttt{xml} (Silesia), \texttt{E.coli} (Canterbury) &
      bpb y tiempos de LZW & RQ3 \\
    \bottomrule
  \end{tabularx}
\end{table}

\subsection{Entorno de ejecución}
\label{sec:entorno}
Todas las mediciones se realizaron en una misma máquina, cuyos datos registra
automáticamente el arnés (\cref{tab:entorno}).

\begin{table}[H]
  \centering
  \caption{Entorno de ejecución (generado por \texttt{BenchmarkRunner}).}
  \label{tab:entorno}
  \small
  \begin{tabular}{@{}l l@{}}
    \toprule
    Elemento & Valor \\
    \midrule
    Procesador & \resCpu \\
    Núcleos lógicos disponibles & \resCores \\
    Máquina virtual Java & \resJvm \\
    Memoria máxima del montículo & \resHeap~MiB \\
    Opciones de la JVM & \texttt{-Xms3g -Xmx3g -XX:+UseParallelGC} \\
    Compilación & Java~17 (\texttt{javac.source}), Apache Ant, JUnit 4.13.2 \\
    Fecha de la campaña & \resDate \\
    \bottomrule
  \end{tabular}
\end{table}

\subsection{Amenazas a la validez}
\label{sec:amenazas}
\begin{itemize}
  \item \emph{Ruido de JIT y GC.} Aunque el calentamiento, las repeticiones y la mediana mitigan
        la variabilidad~\cite{georges2007}, las pausas del recolector o la recompilación pueden
        afectar mediciones individuales; los intervalos de confianza lo hacen visible.
  \item \emph{Una sola máquina.} Los valores absolutos dependen del procesador, la jerarquía de
        caché y la JVM; las conclusiones se basan en tendencias (pendientes, ns/byte relativos)
        más que en tiempos absolutos.
  \item \emph{Datos sintéticos frente a reales.} Las fuentes sintéticas aíslan propiedades,
        pero no representan todos los datos reales; por eso se complementan con texto real y
        con los corpus Canterbury y Silesia.
  \item \emph{Calidad de la implementación.} Los tiempos reflejan estas implementaciones
        concretas (por ejemplo, la decodificación de Huffman bit a bit, sin tablas de búsqueda),
        no el límite de cada algoritmo; las optimizaciones conocidas se discuten en la
        literatura~\cite{moffat1997,moffat2019}.
\end{itemize}

% ---------------------------------------------------------------------
%  5. Resultados — Integrante C
%  Todos los valores provienen de results/ (CSV y macros generados por
%  BenchmarkRunner); build.sh copia ../results en report/results.
% ---------------------------------------------------------------------
\section{Resultados}
\label{sec:resultados}

% Gráfico de E1: tiempo frente a n en escala log-log.
% #1 = archivo CSV, #2 = título, #3 = extremos de la recta de pendiente 1.
\newcommand{\graficoEuno}[3]{%
  \begin{tikzpicture}
    \begin{loglogaxis}[informe, height=0.78\linewidth, log basis x=2,
        xlabel={$n$ [bytes]}, ylabel={Mediana [ms]},
        title={#2}, title style={font=\small\bfseries},
        legend to name=leyendaEuno, legend columns=5,
        legend style={draw=gray!50, /tikz/every even column/.append style={column sep=6pt}}]
      \addplot+[huffenc, error bars/.cd, y dir=plus, y explicit]
        table[col sep=comma, x=n, y=huff_enc_ms, y error=huff_enc_ci] {results/#1};
      \addlegendentry{Huffman codificación}
      \addplot+[huffdec, error bars/.cd, y dir=plus, y explicit]
        table[col sep=comma, x=n, y=huff_dec_ms, y error=huff_dec_ci] {results/#1};
      \addlegendentry{Huffman decodificación}
      \addplot+[lzwenc, error bars/.cd, y dir=plus, y explicit]
        table[col sep=comma, x=n, y=lzw_enc_ms, y error=lzw_enc_ci] {results/#1};
      \addlegendentry{LZW codificación}
      \addplot+[lzwdec, error bars/.cd, y dir=plus, y explicit]
        table[col sep=comma, x=n, y=lzw_dec_ms, y error=lzw_dec_ci] {results/#1};
      \addlegendentry{LZW decodificación}
      \addplot[pendiente] coordinates {#3};
      \addlegendentry{pendiente 1 (referencia)}
    \end{loglogaxis}
  \end{tikzpicture}}

% Gráfico de E1: costo por byte frente a n.
\newcommand{\graficoNspb}[2]{%
  \begin{tikzpicture}
    \begin{semilogxaxis}[informe, height=0.78\linewidth, log basis x=2,
        xlabel={$n$ [bytes]}, ylabel={ns/byte}, ymin=0,
        title={#2}, title style={font=\small\bfseries}, legend pos=north east]
      \addplot[huffenc] table[col sep=comma, x=n, y=huff_enc_nspb] {results/#1};
      \addlegendentry{Huffman cod.}
      \addplot[huffdec] table[col sep=comma, x=n, y=huff_dec_nspb] {results/#1};
      \addlegendentry{Huffman decod.}
      \addplot[lzwenc] table[col sep=comma, x=n, y=lzw_enc_nspb] {results/#1};
      \addlegendentry{LZW cod.}
      \addplot[lzwdec] table[col sep=comma, x=n, y=lzw_dec_nspb] {results/#1};
      \addlegendentry{LZW decod.}
    \end{semilogxaxis}
  \end{tikzpicture}}

% Gráfico de E1: bits por byte frente a n para un algoritmo (#1 = prefijo de columna).
\newcommand{\graficoBpbN}[2]{%
  \begin{tikzpicture}
    \begin{semilogxaxis}[informe, height=0.78\linewidth, log basis x=2,
        xlabel={$n$ [bytes]}, ylabel={bits por byte}, ymin=0,
        title={#2}, title style={font=\small\bfseries},
        legend to name=leyendaBpbN, legend columns=4,
        legend style={draw=gray!50, /tikz/every even column/.append style={column sep=6pt}}]
      \addplot[black!70, mark=*, mark size=1.2pt] table[col sep=comma, x=n, y=#1_bpb] {results/e1_uniform.csv};
      \addlegendentry{\texttt{uniform}}
      \addplot[blue!70!black, mark=square*, mark size=1.2pt] table[col sep=comma, x=n, y=#1_bpb] {results/e1_skewed.csv};
      \addlegendentry{\texttt{skewed}}
      \addplot[green!45!black, mark=triangle*, mark size=1.5pt] table[col sep=comma, x=n, y=#1_bpb] {results/e1_synthtext.csv};
      \addlegendentry{\texttt{synthtext}}
      \addplot[red!70!black, mark=diamond*, mark size=1.5pt] table[col sep=comma, x=n, y=#1_bpb] {results/e1_dickens.csv};
      \addlegendentry{\texttt{dickens}}
    \end{semilogxaxis}
  \end{tikzpicture}}

\subsection{E1: escalamiento del tiempo con el tamaño de la entrada}
\label{sec:res-e1}
\begin{aporte}{C}
La \cref{fig:e1-tiempo} muestra la mediana del tiempo de codificación y de decodificación en
función de $n$ para las cuatro fuentes, en escala log--log; la recta discontinua tiene
pendiente 1 y sirve solo como guía visual de crecimiento lineal (su posición vertical es
arbitraria). Las barras de error indican el semiancho superior del intervalo de confianza del
95\,\%.

\siresultado{e1_dickens.csv}{%
\begin{figure}[H]
  \centering
  \begin{subfigure}[t]{0.49\linewidth}
    \centering
    \graficoEuno{e1_uniform.csv}{\texttt{uniform}}{(16384,0.16384) (16777216,167.77216)}
    \caption{Bytes uniformes i.i.d.}
    \label{fig:e1-uniform}
  \end{subfigure}\hfill
  \begin{subfigure}[t]{0.49\linewidth}
    \centering
    \graficoEuno{e1_skewed.csv}{\texttt{skewed}}{(16384,0.16384) (16777216,167.77216)}
    \caption{Geométrica i.i.d. ($p = 0.15$).}
    \label{fig:e1-skewed}
  \end{subfigure}\\[6pt]
  \begin{subfigure}[t]{0.49\linewidth}
    \centering
    \graficoEuno{e1_synthtext.csv}{\texttt{synthtext}}{(16384,0.16384) (16777216,167.77216)}
    \caption{Texto sintético (Zipf).}
    \label{fig:e1-synthtext}
  \end{subfigure}\hfill
  \begin{subfigure}[t]{0.49\linewidth}
    \centering
    \graficoEuno{e1_dickens.csv}{\texttt{dickens}}{(16384,0.16384) (8388608,83.88608)}
    \caption{Texto real (prefijos de \texttt{dickens}).}
    \label{fig:e1-dickens}
  \end{subfigure}\\[4pt]
  \pgfplotslegendfromname{leyendaEuno}
  \caption{E1: mediana del tiempo frente al tamaño de la entrada (escala log--log).}
  \label{fig:e1-tiempo}
\end{figure}}

% TODO-RESULTS: describir la forma de las curvas de la figura e1-tiempo (paralelismo con la recta de pendiente 1, desviaciones en tamaños pequeños, diferencias entre fuentes y entre codificación/decodificación).
\todoresults{Interpretar la \cref{fig:e1-tiempo}: paralelismo de las curvas con la recta de
pendiente 1, desviaciones en los tamaños pequeños y diferencias entre fuentes y entre
codificación y decodificación.}

La \cref{tab:pendientes} reporta la pendiente del ajuste log--log y su $R^2$ para cada fuente,
algoritmo y operación, calculados con $n \ge 2^{16}$, junto con el costo por byte en el mayor
tamaño medido.

\siresultado{fits.csv}{%
\begin{table}[H]
  \centering
  \caption{E1: pendiente $\hat{k}$ del ajuste $\ln t = \hat{k}\ln n + c$, coeficiente $R^2$ y
           costo por byte en el mayor $n$.}
  \label{tab:pendientes}
  \small
  \pgfplotstabletypeset[
    col sep=comma,
    columns={source,algorithm,operation,slope,r2,nspb_at_max},
    columns/source/.style={string type, column name={Fuente},
      postproc cell content/.append style={/pgfplots/table/@cell content/.add={\ttfamily}{}}},
    columns/algorithm/.style={string type, column name={Algoritmo},
      string replace={huffman}{Huffman}, string replace={lzw}{LZW}},
    columns/operation/.style={string type, column name={Operación},
      string replace={encode}{codificación}, string replace={decode}{decodificación}},
    columns/slope/.style={column name={Pendiente $\hat{k}$}, fixed, fixed zerofill, precision=3},
    columns/r2/.style={column name={$R^2$}, fixed, fixed zerofill, precision=4},
    columns/nspb_at_max/.style={column name={ns/byte en $n_{\max}$}, fixed, fixed zerofill, precision=2},
    every head row/.style={before row=\toprule, after row=\midrule},
    every last row/.style={after row=\bottomrule},
    every row no 3/.style={after row=\midrule},
    every row no 7/.style={after row=\midrule},
    every row no 11/.style={after row=\midrule},
  ]{results/fits.csv}
\end{table}}

% TODO-RESULTS: comentar el rango de pendientes y R^2 de la tabla tab:pendientes (por ejemplo, con \resSlopeDickensHuffEnc, \resSlopeDickensLzwEnc) y su cercanía a 1.
\todoresults{Comentar el rango de pendientes y de $R^2$ de la \cref{tab:pendientes} y su
cercanía a 1 (por ejemplo, en \texttt{dickens}: Huffman \resSlopeDickensHuffEnc{} /
\resSlopeDickensHuffDec{}, LZW \resSlopeDickensLzwEnc{} / \resSlopeDickensLzwDec{}).}

La \cref{fig:e1-nspb} presenta el mismo experimento como costo por byte en escala
semilogarítmica: bajo complejidad lineal, las curvas deben ser aproximadamente horizontales.

\siresultado{e1_synthtext.csv}{%
\begin{figure}[H]
  \centering
  \begin{subfigure}[t]{0.49\linewidth}
    \centering
    \graficoNspb{e1_dickens.csv}{\texttt{dickens}}
    \caption{Texto real.}
    \label{fig:e1-nspb-dickens}
  \end{subfigure}\hfill
  \begin{subfigure}[t]{0.49\linewidth}
    \centering
    \graficoNspb{e1_synthtext.csv}{\texttt{synthtext}}
    \caption{Texto sintético.}
    \label{fig:e1-nspb-synthtext}
  \end{subfigure}
  \caption{E1: costo por byte (ns/byte) frente al tamaño de la entrada.}
  \label{fig:e1-nspb}
\end{figure}}

% TODO-RESULTS: interpretar la figura e1-nspb (estabilidad del costo por byte, efectos de caché o de costos fijos en tamaños pequeños, comparación del costo relativo de ambos algoritmos).
\todoresults{Interpretar la \cref{fig:e1-nspb}: estabilidad del costo por byte, efectos de
costos fijos o de caché y costo relativo de cada algoritmo y operación.}

Finalmente, la \cref{fig:e1-bpb} muestra la tasa de compresión obtenida en E1 para cada
fuente y tamaño.

\siresultado{e1_uniform.csv}{%
\begin{figure}[H]
  \centering
  \begin{subfigure}[t]{0.49\linewidth}
    \centering
    \graficoBpbN{huff}{Huffman}
    \caption{Huffman.}
    \label{fig:e1-bpb-huff}
  \end{subfigure}\hfill
  \begin{subfigure}[t]{0.49\linewidth}
    \centering
    \graficoBpbN{lzw}{LZW ($b = 16$)}
    \caption{LZW.}
    \label{fig:e1-bpb-lzw}
  \end{subfigure}\\[4pt]
  \pgfplotslegendfromname{leyendaBpbN}
  \caption{E1: bits por byte frente al tamaño de la entrada.}
  \label{fig:e1-bpb}
\end{figure}}

% TODO-RESULTS: interpretar la figura e1-bpb (efecto de la cabecera de Huffman en n pequeño, aprendizaje del diccionario de LZW, expansión de LZW en uniform).
\todoresults{Interpretar la \cref{fig:e1-bpb}: efecto de la cabecera de 262 bytes de Huffman en
$n$ pequeño, mejora de LZW a medida que el diccionario aprende y comportamiento en
\texttt{uniform}.}
\end{aporte}

\subsection{E2: barrido de entropía}
\label{sec:res-e2}
\begin{aporte}{C}
En E2 la fuente es i.i.d. uniforme sobre $k$ símbolos, de modo que $H_0 = \log_2 k$ se conoce
por construcción y no hay frases repetidas más allá de las debidas al azar. La
\cref{fig:e2-bpb} compara los bits por byte de ambos algoritmos con la entropía empírica y con
la cota superior $H_0 + 1$ de la \cref{eq:cota-huffman}.

\siresultado{e2_entropy.csv}{%
\begin{figure}[H]
  \centering
  \begin{tikzpicture}
    \begin{semilogxaxis}[informe, width=0.72\linewidth, height=0.5\linewidth, log basis x=2,
        xlabel={$k$ (tamaño del alfabeto); $H_0 = \log_2 k$}, ylabel={bits por byte},
        ymin=0, legend pos=north west]
      \addplot[huffenc] table[col sep=comma, x=k, y=huff_bpb] {results/e2_entropy.csv};
      \addlegendentry{Huffman}
      \addplot[lzwenc] table[col sep=comma, x=k, y=lzw_bpb] {results/e2_entropy.csv};
      \addlegendentry{LZW ($b = 16$)}
      \addplot[black, thick, mark=none] table[col sep=comma, x=k, y=h0] {results/e2_entropy.csv};
      \addlegendentry{$H_0$ empírica}
      \addplot[pendiente] table[col sep=comma, x=k, y expr=\thisrow{h0}+1] {results/e2_entropy.csv};
      \addlegendentry{$H_0 + 1$}
    \end{semilogxaxis}
  \end{tikzpicture}
  \caption{E2: bits por byte frente al tamaño del alfabeto ($n = 2^{23}$).}
  \label{fig:e2-bpb}
\end{figure}}

% TODO-RESULTS: interpretar e2-bpb (Huffman frente a [H0, H0+1), distancia de LZW a H0 por k; citar por ejemplo \resEntropyHuffBpbKTwo, \resEntropyLzwBpbKTwo, \resEntropyHuffBpbKTwoFiftySix, \resEntropyLzwBpbKTwoFiftySix).
\todoresults{Interpretar la \cref{fig:e2-bpb}: posición de Huffman en $[H_0, H_0+1)$ y
distancia de LZW a $H_0$ para cada $k$ (por ejemplo, $k = 2$: Huffman \resEntropyHuffBpbKTwo{},
LZW \resEntropyLzwBpbKTwo{}; $k = 256$: Huffman \resEntropyHuffBpbKTwoFiftySix{}, LZW
\resEntropyLzwBpbKTwoFiftySix{}).}

La \cref{fig:e2-tiempo} muestra cómo varía el tiempo con la entropía para un tamaño fijo, lo
que permite contrastar el término $\Lbar$ del costo de decodificación de Huffman.

\siresultado{e2_entropy.csv}{%
\begin{figure}[H]
  \centering
  \begin{subfigure}[t]{0.49\linewidth}
    \centering
    \begin{tikzpicture}
      \begin{semilogxaxis}[informe, height=0.78\linewidth, log basis x=2,
          xlabel={$k$}, ylabel={Mediana [ms]}, ymin=0, legend pos=north west]
        \addplot+[huffdec, error bars/.cd, y dir=plus, y explicit]
          table[col sep=comma, x=k, y=huff_dec_ms, y error=huff_dec_ci] {results/e2_entropy.csv};
        \addlegendentry{Huffman}
        \addplot+[lzwdec, error bars/.cd, y dir=plus, y explicit]
          table[col sep=comma, x=k, y=lzw_dec_ms, y error=lzw_dec_ci] {results/e2_entropy.csv};
        \addlegendentry{LZW}
      \end{semilogxaxis}
    \end{tikzpicture}
    \caption{Decodificación.}
    \label{fig:e2-dec}
  \end{subfigure}\hfill
  \begin{subfigure}[t]{0.49\linewidth}
    \centering
    \begin{tikzpicture}
      \begin{semilogxaxis}[informe, height=0.78\linewidth, log basis x=2,
          xlabel={$k$}, ylabel={Mediana [ms]}, ymin=0, legend pos=north west]
        \addplot+[huffenc, error bars/.cd, y dir=plus, y explicit]
          table[col sep=comma, x=k, y=huff_enc_ms, y error=huff_enc_ci] {results/e2_entropy.csv};
        \addlegendentry{Huffman}
        \addplot+[lzwenc, error bars/.cd, y dir=plus, y explicit]
          table[col sep=comma, x=k, y=lzw_enc_ms, y error=lzw_enc_ci] {results/e2_entropy.csv};
        \addlegendentry{LZW}
      \end{semilogxaxis}
    \end{tikzpicture}
    \caption{Codificación.}
    \label{fig:e2-enc}
  \end{subfigure}
  \caption{E2: tiempo frente al tamaño del alfabeto ($n = 2^{23}$).}
  \label{fig:e2-tiempo}
\end{figure}}

% TODO-RESULTS: interpretar e2-tiempo (crecimiento del tiempo de decodificación de Huffman con log2 k frente a LZW; citar \resEntropyHuffDecMsKTwo, \resEntropyHuffDecMsKTwoFiftySix, \resEntropyLzwDecMsKTwo, \resEntropyLzwDecMsKTwoFiftySix).
\todoresults{Interpretar la \cref{fig:e2-tiempo}: tendencia del tiempo de decodificación de
Huffman con $\log_2 k$ frente a LZW (por ejemplo, Huffman \resEntropyHuffDecMsKTwo{} ms con
$k = 2$ y \resEntropyHuffDecMsKTwoFiftySix{} ms con $k = 256$; LZW \resEntropyLzwDecMsKTwo{} y
\resEntropyLzwDecMsKTwoFiftySix{} ms).}
\end{aporte}

\subsection{E3: corpus Canterbury y Silesia}
\label{sec:res-e3}
\begin{aporte}{C}
La \cref{tab:corpus} reporta, para cada archivo, su tamaño, su entropía de orden cero, los bits
por byte de cada algoritmo y el rendimiento de codificación. En negrita aparece el mejor de los
dos algoritmos comparados; DEFLATE se muestra solo como referencia.

\siresultado{table_corpus.tex}{%
\begin{table}[H]
  \centering
  \caption{E3: bits por byte y rendimiento de codificación en los corpus Canterbury (arriba) y
           Silesia (abajo).}
  \label{tab:corpus}
  \small
  \setlength{\tabcolsep}{5pt}
  \begin{tabular}{@{}l r r r r r r r@{}}
    \toprule
    Archivo & Tamaño (KiB) & $H_0$ & Huffman & LZW & DEFLATE$^\dagger$ & Huff. MB/s & LZW MB/s \\
    \midrule
    \texttt{E.coli} & 4530 & 2.00 & \textbf{2.00} & 2.17 & 2.31 & 348 & 141 \\
\texttt{alice29.txt} & 149 & 4.57 & 4.63 & \textbf{3.27} & 2.86 & 222 & 105 \\
\texttt{asyoulik.txt} & 122 & 4.81 & 4.86 & \textbf{3.51} & 3.12 & 228 & 107 \\
\texttt{bible.txt} & 3953 & 4.34 & 4.39 & \textbf{2.78} & 2.36 & 234 & 98 \\
\texttt{cp.html} & 24 & 5.23 & 5.35 & \textbf{3.68} & 2.59 & 231 & 89 \\
\texttt{fields.c} & 11 & 5.01 & 5.23 & \textbf{3.56} & 2.24 & 206 & 69 \\
\texttt{grammar.lsp} & 4 & 4.63 & 5.23 & \textbf{3.91} & 2.62 & 214 & 32 \\
\texttt{kennedy.xls} & 1006 & 3.57 & 3.60 & \textbf{2.40} & 1.58 & 286 & 142 \\
\texttt{lcet10.txt} & 417 & 4.67 & 4.70 & \textbf{3.19} & 2.72 & 235 & 102 \\
\texttt{plrabn12.txt} & 471 & 4.53 & 4.58 & \textbf{3.41} & 3.24 & 228 & 104 \\
\texttt{ptt5} & 501 & 1.21 & 1.66 & \textbf{0.97} & 0.88 & 251 & 203 \\
\texttt{sum} & 37 & 5.33 & 5.42 & \textbf{4.21} & 2.72 & 186 & 91 \\
\texttt{world192.txt} & 2415 & 5.00 & 5.04 & \textbf{3.26} & 2.34 & 196 & 96 \\
\texttt{xargs.1} & 4 & 4.90 & 5.42 & \textbf{4.43} & 3.28 & 91 & 5 \\
\midrule
\texttt{dickens} & 9954 & 4.53 & 4.57 & \textbf{3.29} & 3.04 & 216 & 95 \\
\texttt{mozilla} & 50020 & 6.22 & 6.24 & \textbf{4.30} & 2.98 & 150 & 63 \\
\texttt{mr} & 9737 & 3.68 & 3.71 & \textbf{3.01} & 2.95 & 246 & 112 \\
\texttt{nci} & 32767 & 2.43 & 2.44 & \textbf{0.94} & 0.76 & 264 & 137 \\
\texttt{ooffice} & 6008 & 6.64 & 6.66 & \textbf{5.29} & 4.03 & 205 & 83 \\
\texttt{osdb} & 9849 & 6.59 & 6.62 & \textbf{4.38} & 2.93 & 219 & 84 \\
\texttt{reymont} & 6472 & 4.84 & 4.87 & \textbf{2.46} & 2.25 & 218 & 106 \\
\texttt{samba} & 21100 & 6.09 & 6.13 & \textbf{3.06} & 2.02 & 212 & 100 \\
\texttt{sao} & 7082 & 7.53 & 7.55 & \textbf{7.53} & 5.88 & 214 & 80 \\
\texttt{webster} & 40487 & 4.97 & 5.00 & \textbf{2.76} & 2.36 & 207 & 95 \\
\texttt{x-ray} & 8276 & 6.60 & \textbf{6.63} & 6.67 & 5.71 & 237 & 102 \\
\texttt{xml} & 5220 & 5.52 & 5.55 & \textbf{1.66} & 1.03 & 247 & 103 \\

    \bottomrule
  \end{tabular}
  \par\smallskip
  \begin{minipage}{0.92\linewidth}
    \footnotesize $^\dagger$\,DEFLATE (zlib del JDK, nivel 6, LZ77 + Huffman) es solo una
    referencia externa y no forma parte de la comparación. Las columnas Huffman, LZW y DEFLATE
    están en bits por byte, incluida la cabecera.
  \end{minipage}
\end{table}}

En conjunto, sobre los \resCorpusFiles{} archivos, la tasa agregada es de
\resCorpusHuffBpb{}~bpb para Huffman, \resCorpusLzwBpb{}~bpb para LZW y
\resCorpusDeflateBpb{}~bpb para DEFLATE; LZW comprime mejor que Huffman en \resLzwWins{} de
los \resCorpusFiles{} archivos.

% TODO-RESULTS: interpretar tab:corpus (en qué tipos de archivo gana cada algoritmo, relación con H0, archivos donde LZW expande o se acerca a Huffman, rendimiento relativo en MB/s, distancia a DEFLATE).
\todoresults{Interpretar la \cref{tab:corpus}: tipos de archivo en que gana cada algoritmo,
relación con $H_0$, casos en que LZW se acerca o supera a 8~bpb, rendimiento relativo en MB/s
y distancia de ambos a DEFLATE.}
\end{aporte}

\subsection{E4: tamaño máximo del diccionario de LZW}
\label{sec:res-e4}
\begin{aporte}{C}
La \cref{fig:e4} muestra los bits por byte y el tiempo de codificación de LZW al variar el
ancho máximo de código $b$, es decir, el tamaño del diccionario $2^b$, en tres archivos de
naturaleza distinta. El ancho con menor bpb fue $b = \resLzwBitsBestDickens$ en
\texttt{dickens}, $b = \resLzwBitsBestXml$ en \texttt{xml} y $b = \resLzwBitsBestEcoli$ en
\texttt{E.coli}.

\siresultado{e4_lzw_bits.csv}{%
\begin{figure}[H]
  \centering
  \begin{subfigure}[t]{0.49\linewidth}
    \centering
    \begin{tikzpicture}
      \begin{axis}[informe, height=0.78\linewidth, xlabel={$b$ (bits; diccionario de $2^b$ entradas)},
          ylabel={bits por byte}, xtick={9,...,16}, legend pos=north east]
        \addplot[blue!70!black, mark=*, mark size=1.3pt]
          table[col sep=comma, x=bits, y=dickens_bpb] {results/e4_lzw_bits.csv};
        \addlegendentry{\texttt{dickens}}
        \addplot[green!45!black, mark=square*, mark size=1.2pt]
          table[col sep=comma, x=bits, y=xml_bpb] {results/e4_lzw_bits.csv};
        \addlegendentry{\texttt{xml}}
        \addplot[red!70!black, mark=triangle*, mark size=1.5pt]
          table[col sep=comma, x=bits, y=ecoli_bpb] {results/e4_lzw_bits.csv};
        \addlegendentry{\texttt{E.coli}}
      \end{axis}
    \end{tikzpicture}
    \caption{Tasa de compresión.}
    \label{fig:e4-bpb}
  \end{subfigure}\hfill
  \begin{subfigure}[t]{0.49\linewidth}
    \centering
    \begin{tikzpicture}
      \begin{axis}[informe, height=0.78\linewidth, xlabel={$b$ (bits; diccionario de $2^b$ entradas)},
          ylabel={Mediana de codificación [ms]}, xtick={9,...,16}, ymin=0, legend pos=north east]
        \addplot[blue!70!black, mark=*, mark size=1.3pt]
          table[col sep=comma, x=bits, y=dickens_enc_ms] {results/e4_lzw_bits.csv};
        \addlegendentry{\texttt{dickens}}
        \addplot[green!45!black, mark=square*, mark size=1.2pt]
          table[col sep=comma, x=bits, y=xml_enc_ms] {results/e4_lzw_bits.csv};
        \addlegendentry{\texttt{xml}}
        \addplot[red!70!black, mark=triangle*, mark size=1.5pt]
          table[col sep=comma, x=bits, y=ecoli_enc_ms] {results/e4_lzw_bits.csv};
        \addlegendentry{\texttt{E.coli}}
      \end{axis}
    \end{tikzpicture}
    \caption{Tiempo de codificación.}
    \label{fig:e4-tiempo}
  \end{subfigure}
  \caption{E4: efecto del tamaño máximo del diccionario de LZW.}
  \label{fig:e4}
\end{figure}}

% TODO-RESULTS: interpretar fig:e4 (tendencia del bpb con b, rendimientos decrecientes, diferencias entre texto, XML y ADN, costo temporal de diccionarios grandes y frecuencia de reinicios).
\todoresults{Interpretar la \cref{fig:e4}: tendencia de los bits por byte al crecer $b$,
rendimientos decrecientes, diferencias entre texto, XML y ADN, y costo temporal asociado al
tamaño del diccionario y a la frecuencia de reinicios.}
\end{aporte}

% ---------------------------------------------------------------------
%  6. Discusión — C (H1), A (H2, H4), B (H3), D (síntesis por pregunta)
% ---------------------------------------------------------------------
\section{Discusión: teoría frente a experimento}
\label{sec:discusion}

Esta sección contrasta cada hipótesis de la \cref{sec:hipotesis} con la evidencia de la
\cref{sec:resultados} y, a continuación, sintetiza la respuesta a cada pregunta de
investigación.

\subsection{H1: escalamiento lineal}
\label{sec:disc-h1}
\begin{aporte}{C}
\paragraph{Predicción.} Según la \cref{tab:comparacion-teorica}, los cuatro procedimientos son
lineales en $n$: la codificación de Huffman es $\Theta(n + \sigma\log\sigma)$, su
decodificación $\Theta(n\Lbar)$ con $\Lbar$ acotado, y ambas operaciones de LZW son
$\Theta(n)$. En escala log--log, la pendiente esperada es 1.

\paragraph{Desviaciones previsibles.} El análisis asintótico ignora costos fijos que pesan en
entradas pequeñas: Huffman construye el código y escribe una cabecera de 262 bytes en cada
llamada, y LZW reserva e inicializa en cada llamada su tabla hash de $2^{17}$ posiciones (dos
arreglos de enteros) o, al decodificar, cuatro arreglos de $2^{16}$ entradas. Estos costos
constantes aplanan la curva en los primeros tamaños y tienden a producir pendientes menores
que 1 si se incluyen en el ajuste; por eso el ajuste usa $n \ge 2^{16}$. En el extremo opuesto,
cuando los datos y las tablas dejan de caber en la memoria caché, el costo por byte puede
aumentar ligeramente sin que cambie la complejidad.

\paragraph{Evidencia.} \Cref{fig:e1-tiempo,fig:e1-nspb} y \cref{tab:pendientes}.
% TODO-RESULTS: contrastar H1 con las pendientes y R^2 medidos (tab:pendientes), indicar si se acepta, y explicar cualquier pendiente alejada de 1 con los costos fijos o efectos de caché descritos.
\todoresults{Contrastar H1 con las pendientes y $R^2$ de la \cref{tab:pendientes}; indicar si
se acepta y explicar cualquier pendiente alejada de 1 con los costos fijos o los efectos de
caché descritos.}
\end{aporte}

\subsection{H2: Huffman frente a la entropía}
\label{sec:disc-h2}
\begin{aporte}{A}
\paragraph{Predicción.} La tasa medida incluye la cabecera y el relleno del último byte, por lo
que
\begin{equation}
  \bpb_{\text{Huffman}} = \Lbar + \frac{8 \cdot 262}{n} + \frac{\delta}{n}, \qquad 0 \le \delta < 8,
  \label{eq:bpb-huffman}
\end{equation}
y, por la \cref{eq:cota-huffman}, tiende a un valor en $[H_0, H_0 + 1)$ cuando $n$ crece. En
E2 con $k$ potencia de dos y $n = 2^{23}$, las frecuencias empíricas son casi iguales, el árbol
resultante es completo y equilibrado, y cada símbolo recibe $\log_2 k$ bits: la teoría predice
una tasa prácticamente igual a $H_0$, con un sobrecosto de cabecera de $8\cdot 262/2^{23}$ bpb,
del orden de $2.5\times10^{-4}$. En la fuente geométrica, en cambio, las probabilidades no son
potencias de $1/2$ y cabe esperar una distancia positiva, aunque menor que un bit, respecto de
$H_0$. Sobre texto, Huffman no puede bajar de $H_0$ porque no modela el contexto.

\paragraph{Evidencia.} \Cref{fig:e2-bpb,fig:e1-bpb} y la columna $H_0$ de la \cref{tab:corpus}.
% TODO-RESULTS: verificar que todos los bpb de Huffman (E1, E2, E3) estén en [H0, H0+1) salvo el efecto de cabecera en n pequeño; cuantificar la distancia a H0 en E2 y en skewed.
\todoresults{Verificar que los bits por byte de Huffman en E1, E2 y E3 queden en
$[H_0, H_0+1)$, salvo el efecto de la cabecera en $n$ pequeño, y cuantificar la distancia a
$H_0$ en E2 y en \texttt{skewed}.}
\end{aporte}

\subsection{H3: tipo de redundancia y ganador}
\label{sec:disc-h3}
\begin{aporte}{B}
\paragraph{Predicción.} LZW solo ahorra espacio cuando sus frases son largas: cada código
cuesta entre 9 y $b$ bits, así que una frase de longitud media $m$ cuesta del orden de
$b/m$ bits por byte. En datos con palabras y frases recurrentes, el diccionario captura
secuencias largas y la tasa puede quedar por debajo de $H_0$, algo imposible para un código de
orden cero. En datos i.i.d. las frases crecen muy lentamente, porque cada extensión requiere
que un par (frase, byte) ya haya aparecido; con bytes uniformes, la mayoría de las frases
tienen uno o dos bytes y cada código de 9 a 16 bits representa poco más de un byte, de modo
que se espera expansión ($\bpb > 8$). Huffman, por el contrario, nunca supera
$\Lbar < H_0 + 1$ más la cabecera.

\paragraph{Evidencia.} \Cref{fig:e1-bpb,fig:e2-bpb} y \cref{tab:corpus}.
% TODO-RESULTS: contrastar H3 con E1 (synthtext, dickens frente a uniform, skewed), E2 y E3 (\resLzwWins de \resCorpusFiles); identificar archivos del corpus que contradigan la hipótesis y explicar por qué.
\todoresults{Contrastar H3 con E1 (\texttt{synthtext} y \texttt{dickens} frente a
\texttt{uniform} y \texttt{skewed}), con E2 y con E3 (LZW gana en \resLzwWins{} de
\resCorpusFiles{} archivos); identificar los archivos que contradigan la hipótesis y explicar
por qué.}
\end{aporte}

\subsection{H4: sensibilidad del tiempo a la entropía}
\label{sec:disc-h4}
\begin{aporte}{A}
\paragraph{Predicción.} El decodificador canónico ejecuta una iteración por bit leído, $n\Lbar$
en total; con $\Lbar \approx \log_2 k$ en E2, su tiempo debería crecer aproximadamente de forma
lineal con $\log_2 k$. El codificador de Huffman escribe un campo por símbolo con
\texttt{BitWriter}, cuyo costo depende solo débilmente de $\ell_s$, por lo que se espera
menor sensibilidad. La decodificación de LZW escribe cada byte de salida una vez; la entropía
solo afecta al número de códigos, que crece cuando las frases se acortan, de modo que se espera
una dependencia mucho más débil que la de Huffman.

\paragraph{Evidencia.} \Cref{fig:e2-tiempo}.
% TODO-RESULTS: contrastar H4 con E2 (pendiente aproximada del tiempo de decodificación de Huffman frente a log2 k; comparar con LZW usando \resEntropyHuffDecMsKTwo, \resEntropyHuffDecMsKTwoFiftySix, \resEntropyLzwDecMsKTwo, \resEntropyLzwDecMsKTwoFiftySix).
\todoresults{Contrastar H4 con la \cref{fig:e2-tiempo}: tendencia del tiempo de decodificación
de Huffman frente a $\log_2 k$ comparada con la de LZW.}
\end{aporte}

\subsection{Síntesis por pregunta de investigación}
\label{sec:disc-rq}
\begin{aporte}{D}
\paragraph{RQ1 (escalamiento).} Se responde con H1 y H4: la teoría predice crecimiento lineal
para ambos algoritmos, con una constante que en la decodificación de Huffman depende de la
longitud media del código.
% TODO-RESULTS: respuesta a RQ1 basada en H1 y H4.
\todoresults{Responder RQ1 a partir de la discusión de H1 y H4.}

\paragraph{RQ2 (compresión, entropía y tipo de dato).} Se responde con H2 y H3: Huffman está
acotado por $H_0$ y no puede aprovechar repeticiones, mientras que LZW puede bajar de $H_0$ en
datos con frases repetidas, pero pierde ---y puede expandir--- en fuentes sin memoria.
% TODO-RESULTS: respuesta a RQ2 basada en H2, H3 y E3.
\todoresults{Responder RQ2 a partir de H2, H3 y del agregado del corpus.}

\paragraph{RQ3 (tamaño del diccionario).} Un diccionario mayor retiene más frases y reinicia con
menos frecuencia, pero cada código cuesta más bits; la teoría no fija un óptimo universal, que
depende de cuán repetitivos y estacionarios sean los datos~\cite{cheng2021}.
% TODO-RESULTS: respuesta a RQ3 basada en E4 (\resLzwBitsBestDickens, \resLzwBitsBestXml, \resLzwBitsBestEcoli).
\todoresults{Responder RQ3 a partir de E4 (mejor ancho: \texttt{dickens}
$b = \resLzwBitsBestDickens$, \texttt{xml} $b = \resLzwBitsBestXml$, \texttt{E.coli}
$b = \resLzwBitsBestEcoli$).}
\end{aporte}

% ---------------------------------------------------------------------
%  7. Conclusiones — Integrante D
% ---------------------------------------------------------------------
\section{Conclusiones}
\label{sec:conclusiones}

\subsection{Comparación entre el análisis teórico y los resultados experimentales}
\label{sec:concl-teoria-experimento}
\begin{aporte}{D}
La \cref{tab:teoria-experimento} reúne, para cada aspecto analizado, la predicción derivada en
la \cref{sec:marco-teorico} y la evidencia experimental correspondiente.

\begin{table}[H]
  \centering
  \caption{Predicción teórica frente a evidencia experimental.}
  \label{tab:teoria-experimento}
  \small
  \renewcommand{\arraystretch}{1.15}
  \begin{tabular}{@{}>{\raggedright}p{3.0cm} >{\raggedright}p{5.2cm} >{\raggedright\arraybackslash}p{6.4cm}@{}}
    \toprule
    Aspecto & Predicción teórica & Evidencia experimental \\
    \midrule
    Tiempo de Huffman & $\Theta(n)$ al codificar; $\Theta(n\Lbar)$ al decodificar (H1) &
      Pendiente en \texttt{dickens}: \resSlopeDickensHuffEnc{} (cod.), \resSlopeDickensHuffDec{}
      (decod.); $R^2$: \resRtwoDickensHuffEnc{}, \resRtwoDickensHuffDec{}; resto en la
      \cref{tab:pendientes} \\
    Tiempo de LZW & $\Theta(n)$ esperado al codificar; $\Theta(n)$ al decodificar (H1) &
      Pendiente en \texttt{dickens}: \resSlopeDickensLzwEnc{} (cod.), \resSlopeDickensLzwDec{}
      (decod.); $R^2$: \resRtwoDickensLzwEnc{}, \resRtwoDickensLzwDec{} \\
    Costo por byte & Constante en $n$ & En el mayor prefijo de \texttt{dickens}: Huffman
      \resNspbDickensHuffEnc{} / \resNspbDickensHuffDec{} ns/byte; LZW
      \resNspbDickensLzwEnc{} / \resNspbDickensLzwDec{} ns/byte (cod. / decod.) \\
    Tasa de Huffman & $[H_0, H_0 + 1)$; cerca de $H_0$ en fuentes i.i.d. (H2) &
      E2: \resEntropyHuffBpbKTwo{} bpb con $k = 2$ y \resEntropyHuffBpbKTwoFiftySix{} bpb con
      $k = 256$ \\
    Tasa de LZW según el dato & Mejor con frases repetidas; peor en i.i.d.; expansión en
      bytes uniformes (H3) &
      E2: \resEntropyLzwBpbKTwo{} bpb con $k = 2$ y \resEntropyLzwBpbKTwoFiftySix{} bpb con
      $k = 256$; E3: gana en \resLzwWins{} de \resCorpusFiles{} archivos \\
    Corpus completo & Sin predicción universal &
      Agregado: Huffman \resCorpusHuffBpb{}, LZW \resCorpusLzwBpb{}, DEFLATE
      \resCorpusDeflateBpb{} bpb \\
    Decodificación y entropía & Huffman crece con $\Lbar$; LZW casi insensible (H4) &
      E2, decodificación: Huffman \resEntropyHuffDecMsKTwo{} a
      \resEntropyHuffDecMsKTwoFiftySix{} ms; LZW \resEntropyLzwDecMsKTwo{} a
      \resEntropyLzwDecMsKTwoFiftySix{} ms ($k = 2$ a $256$) \\
    \bottomrule
  \end{tabular}
\end{table}

% TODO-RESULTS: redactar el párrafo de contraste teoría-experimento: qué predicciones se confirmaron, cuáles se cumplieron parcialmente y por qué (costos fijos, caché, implementación bit a bit).
\todoresults{Redactar el contraste entre teoría y experimento a partir de la
\cref{tab:teoria-experimento}: predicciones confirmadas, cumplidas parcialmente y sus causas.}
\end{aporte}

\subsection{Respuestas a las preguntas de investigación}
\label{sec:concl-rq}
\begin{aporte}{D}
\begin{enumerate}[label=\textbf{RQ\arabic*.}, leftmargin=*]
  \item % TODO-RESULTS: conclusión sobre el escalamiento temporal de ambos algoritmos.
        \todoresults{Conclusión sobre el escalamiento temporal de ambos algoritmos frente a su
        complejidad teórica.}
  \item % TODO-RESULTS: conclusión sobre la relación entre tasa de compresión, entropía y tipo de dato.
        \todoresults{Conclusión sobre la relación entre tasa de compresión, entropía de orden
        cero y tipo de dato; recomendación de uso de cada algoritmo.}
  \item % TODO-RESULTS: conclusión sobre el tamaño del diccionario de LZW.
        \todoresults{Conclusión sobre el efecto del tamaño máximo del diccionario de LZW.}
\end{enumerate}

En términos generales, el análisis teórico indica que la elección entre Huffman y LZW no
depende de su complejidad asintótica, que es lineal en ambos casos, sino de la coincidencia
entre la redundancia que cada uno explota y la que contienen los datos, en línea con la
observación de las revisiones del área~\cite{bell1989,jayasankar2021}.
% TODO-RESULTS: confirmar o matizar esta conclusión general a la luz de los resultados.
\todoresults{Confirmar o matizar la conclusión general a la luz de los resultados.}
\end{aporte}

\subsection{Limitaciones}
\label{sec:limitaciones}
\begin{aporte}{D}
\begin{itemize}
  \item Las mediciones provienen de una sola máquina y una sola JVM; los tiempos absolutos no
        son transferibles, aunque las tendencias sí deberían serlo.
  \item Las implementaciones priorizan la claridad: la decodificación de Huffman lee un bit por
        iteración en lugar de usar tablas de búsqueda, y LZW usa solo la política de reinicio.
  \item Huffman se aplicó con un modelo de orden cero sobre bytes; no se evaluaron modelos con
        contexto ni alfabetos distintos del byte.
  \item DEFLATE se usó solo como referencia de tasa, mediante una implementación nativa no
        comparable en tiempo.
\end{itemize}
\end{aporte}

\subsection{Trabajo futuro}
\label{sec:trabajo-futuro}
\begin{aporte}{D}
\begin{itemize}
  \item Implementar Huffman adaptativo~\cite{vitter1987} para eliminar la segunda pasada y la
        cabecera, y comparar su costo por símbolo con la versión estática.
  \item Incorporar codificación aritmética~\cite{witten1987} para medir cuánto de la distancia
        entre Huffman y $H_0$ se debe a la restricción de bits enteros.
  \item Comparar con LZ77 y con una implementación propia de DEFLATE~\cite{ziv1977,rfc1951},
        que combina diccionario y codificación estadística.
  \item Evaluar políticas alternativas de diccionario lleno en LZW, como LRU o
        IDA-LZW~\cite{cheng2021}, frente al reinicio usado aquí.
  \item Acelerar la decodificación canónica con tablas de búsqueda~\cite{moffat1997} y repetir
        E2 para cuantificar su efecto sobre H4.
\end{itemize}
\end{aporte}


% ---------------------------------------------------------------- bibliografía
\clearpage
\phantomsection
\addcontentsline{toc}{section}{Referencias}
\bibliographystyle{IEEEtran}
\bibliography{refs}

% ---------------------------------------------------------------- anexos
\clearpage
\appendix
% ---------------------------------------------------------------------
%  Anexos — A: reproducibilidad (C); B: aportes individuales (A, B, C, D)
% ---------------------------------------------------------------------
\section{Reproducibilidad}
\label{app:reproducibilidad}
\begin{aporte}{C}
El repositorio contiene todo lo necesario para repetir el estudio. Desde la raíz, con
\texttt{bash}, la secuencia completa es la del \cref{lst:reproducir}. El archivo
\texttt{env.sh} configura el JDK, Apache Ant y Tectonic incluidos en \texttt{.tools/}; el
\emph{script} de corpus descarga Canterbury, el corpus grande y Silesia, y verifica sus sumas
SHA-256; \texttt{ant benchmark} ejecuta E1--E4 y escribe en \texttt{results/} los CSV, la tabla
del corpus y las macros que este informe importa; \texttt{build.sh} copia esos resultados y
compila el PDF.

\begin{lstlisting}[language=bash, numbers=none, caption={Comandos para reproducir el estudio.}, label={lst:reproducir}]
source env.sh                  # JDK, Ant y Tectonic portables
data/download-corpora.sh       # corpus Canterbury, grande y Silesia (SHA-256)
cd code
ant test                       # 43 pruebas JUnit
ant benchmark                  # experimentos E1-E4 -> ../results
cd ../report
./build.sh                     # copia ../results y genera build/main.pdf
\end{lstlisting}

Para una verificación rápida del arnés, sin valor estadístico, puede ejecutarse
\texttt{ant benchmark -Dbench.args="--quick"}, y la opción \texttt{--only e1,e2} restringe los
experimentos. El proyecto también se abre directamente en NetBeans (\emph{File > Open
Project}, carpeta \texttt{code/}), donde \emph{Run} ejecuta una demostración de ambos
algoritmos y \emph{Test} ejecuta la batería de pruebas. La interfaz de línea de comandos
permite usar los algoritmos sobre archivos arbitrarios:
\texttt{java -jar dist/HuffmanVsLZW.jar compress <huffman|lzw|lzw-N> <entrada> <salida>},
\texttt{decompress <entrada> <salida>} y \texttt{stats <entrada>}.
\end{aporte}

\section{Aportes individuales detallados}
\label{app:aportes}
Cada integrante describe a continuación la implementación que desarrolló: su objetivo, su
funcionamiento, cómo se implementó y cómo se integra con el resto del trabajo.

\subsection{\texorpdfstring{\integranteA}{Integrante A}: codificación de Huffman}
\label{app:aporte-a}
\begin{aporte}{A}
\paragraph{Objetivo.} Implementar un compresor de Huffman estático y canónico sobre el alfabeto
de bytes que sea correcto para cualquier entrada, que alcance la longitud media óptima
(\cref{eq:cota-huffman}) y cuya cabecera sea compacta y verificable.

\paragraph{Cómo funciona.} La codificación hace dos pasadas: la primera cuenta las frecuencias
de los 256 bytes; con ellas se calculan las longitudes óptimas mediante el algoritmo voraz
(\cref{alg:huffman}) y se asignan códigos canónicos (\cref{eq:canonico}); la segunda pasada
escribe el código de cada byte. La decodificación lee las 256 longitudes de la cabecera,
reconstruye las tablas canónicas y decodifica bit a bit (\cref{alg:huffman-dec}).

\paragraph{Cómo se implementó.} \texttt{HuffmanCodec} usa un \texttt{PriorityQueue} de índices
de nodo con desempate por índice, arreglos planos de pesos y padres (sin objetos nodo), y un
recorrido descendente para las profundidades (\cref{lst:huffman}). El formato del flujo es:
magia \texttt{HF}, longitud original de 4 bytes y 256 longitudes de 1 byte, seguidas de los
bits. La tabla canónica rechaza longitudes sobresuscritas (suma de Kraft mayor que 1), lo que
detecta cabeceras corruptas. \texttt{BitWriter} y \texttt{BitReader} manejan campos de 1 a 64
bits con un acumulador de 64 bits, bit más significativo primero.

\paragraph{Integración.} \texttt{HuffmanCodec} implementa la interfaz \texttt{Codec}, de modo que
el arnés (\texttt{BenchmarkRunner}) y la línea de comandos (\texttt{Main}) lo usan sin conocer
sus detalles. \texttt{BitWriter} y \texttt{BitReader} son compartidos con LZW. Sus pruebas
(\texttt{HuffmanCodecTest}, 11; \texttt{BitStreamTest}, 5) forman parte de \texttt{ant test}.

\paragraph{Archivos.} \texttt{code/src/pe/edu/ulima/compression/huffman/HuffmanCodec.java},
\texttt{code/src/pe/edu/ulima/compression/io/BitWriter.java},
\texttt{code/src/pe/edu/ulima/compression/io/BitReader.java} y sus pruebas en
\texttt{code/test/pe/edu/ulima/compression/huffman/} y \texttt{.../io/}; en el informe,
\cref{sec:huffman}.
\end{aporte}

\subsection{\texorpdfstring{\integranteB}{Integrante B}: algoritmo LZW}
\label{app:aporte-b}
\begin{aporte}{B}
\paragraph{Objetivo.} Implementar LZW con códigos de ancho variable y tamaño máximo de
diccionario configurable, de forma que la memoria quede acotada y que el ancho máximo pueda
variarse en el experimento E4.

\paragraph{Cómo funciona.} El codificador extiende la frase actual mientras el par (frase,
byte) esté en el diccionario y, cuando no lo está, emite el código de la frase y registra la
nueva entrada (\cref{alg:lzw-enc}). Los códigos crecen de 9 a $b$ bits; al llenarse el
diccionario se emite \texttt{CLEAR} (código 256) y se reinicia. El decodificador reconstruye el
mismo diccionario con un código de retraso, resuelve el caso $KwKwK$ y escribe cada frase de
atrás hacia adelante siguiendo los prefijos (\cref{alg:lzw-dec}).

\paragraph{Cómo se implementó.} \texttt{LzwCodec} representa el diccionario del codificador con
\texttt{PhraseTable}, una tabla hash de direccionamiento abierto de capacidad $2^{b+1}$ (factor de
carga máximo 0.5), sondeo lineal y hashing de Fibonacci; la búsqueda devuelve la ranura donde
insertar si la clave no existe (\cref{lst:lzw}). El decodificador usa cuatro arreglos indexados
por código (\texttt{prefix}, \texttt{suffix}, \texttt{firstByte}, \texttt{phraseLength}). El
formato es: magia \texttt{LZ}, $b$ (1 byte) y longitud original (4 bytes), seguidos de los
códigos. El decodificador valida cada código frente al tamaño actual del diccionario y la
longitud declarada, y rechaza flujos corruptos. El método \texttt{traceCodes} expone la
secuencia emitida para pruebas y enseñanza.

\paragraph{Integración.} \texttt{LzwCodec} implementa \texttt{Codec} y reutiliza
\texttt{BitWriter} y \texttt{BitReader}. El arnés lo instancia con $b = 16$ en E1--E3 y con
$b = 9,\dots,16$ en E4; la línea de comandos acepta \texttt{lzw} o \texttt{lzw-N}. Sus 8 pruebas
(\texttt{LzwCodecTest}) forman parte de \texttt{ant test}.

\paragraph{Archivos.} \texttt{code/src/pe/edu/ulima/compression/lzw/LzwCodec.java} y
\texttt{code/test/pe/edu/ulima/compression/lzw/LzwCodecTest.java}; en el informe,
\cref{sec:lzw}.
\end{aporte}

\subsection{\texorpdfstring{\integranteC}{Integrante C}: metodología experimental y arnés de medición}
\label{app:aporte-c}
\begin{aporte}{C}
\paragraph{Objetivo.} Medir ambos algoritmos de forma estadísticamente rigurosa y
reproducible, y producir automáticamente los datos, tablas y figuras del informe.

\paragraph{Cómo funciona.} \texttt{BenchmarkRunner} ejecuta E1--E4 según el protocolo de la
\cref{sec:protocolo}: calentamiento global, verificación de ida y vuelta, calentamiento local,
repeticiones cronometradas y resúmenes estadísticos. \texttt{Stats} calcula media, mediana,
desviación estándar, intervalo de confianza del 95\,\% con la $t$ de Student, el ajuste
log--log por mínimos cuadrados y la entropía de orden cero. \texttt{DataGenerator} produce las
fuentes sintéticas deterministas de la \cref{tab:datos}.

\paragraph{Cómo se implementó.} El arnés trabaja sobre la interfaz \texttt{Codec}, lo que
permite medir igual a Huffman, a LZW y a DEFLATE (\texttt{DeflateCodec}, envoltorio de
\texttt{java.util.zip}). \texttt{ResultWriter} escribe los CSV, la tabla del corpus en \LaTeX{} y
\texttt{macros.tex}, de modo que ningún número se transcribe a mano. \texttt{Main} ofrece la
línea de comandos con códigos de salida 0, 1 y 2. El \emph{script}
\texttt{data/download-corpora.sh} descarga y verifica los corpus.

\paragraph{Integración.} El arnés consume los códecs de los integrantes A y B; sus salidas en
\texttt{results/} alimentan directamente las figuras y tablas de la \cref{sec:resultados}
mediante \texttt{pgfplots}, \texttt{pgfplotstable} y las macros, a través de
\texttt{report/build.sh}. Sus pruebas (\texttt{StatsTest}, 6; \texttt{DataGeneratorTest}, 3;
\texttt{BenchmarkRunnerTest}, 4; \texttt{MainTest}, 6) forman parte de \texttt{ant test}.

\paragraph{Archivos.} \texttt{code/src/pe/edu/ulima/compression/bench/} (\texttt{BenchmarkRunner},
\texttt{Stats}, \texttt{DataGenerator}, \texttt{DeflateCodec}, \texttt{ResultWriter}),
\texttt{code/src/pe/edu/ulima/compression/app/Main.java}, sus pruebas,
\texttt{data/download-corpora.sh} y \texttt{report/build.sh}; en el informe,
\cref{sec:comparacion-teorica,sec:hipotesis,sec:metodologia,sec:resultados} y la discusión de H1.
\end{aporte}

\subsection{\texorpdfstring{\integranteD}{Integrante D}: estado del arte, redacción y bibliografía}
\label{app:aporte-d}
\begin{aporte}{D}
\paragraph{Objetivo.} Situar el trabajo en la literatura, plantear el problema, los objetivos y
las preguntas de investigación, y redactar las conclusiones a partir del contraste entre teoría
y experimento.

\paragraph{Cómo funciona.} La revisión se organiza en tres niveles ---el problema, los
algoritmos que lo resuelven y las comparaciones previas--- y concluye con la brecha que este
trabajo atiende (\cref{sec:estado-arte}). Las conclusiones responden cada pregunta de
investigación con la evidencia de las \cref{sec:resultados,sec:discusion}.

\paragraph{Cómo se implementó.} La bibliografía (\texttt{report/refs.bib}) reúne 28 referencias
verificadas por DOI o URL: artículos de revista y de congreso, una tesis doctoral, una tesis de
maestría, un RFC y libros de texto; el estilo numérico IEEE se aplica con BibTeX. La estructura
del documento y el código de colores se definen en \texttt{report/main.tex}.

\paragraph{Integración.} Las preguntas de investigación definidas en la introducción
determinan las hipótesis H1--H4 y los experimentos E1--E4; la síntesis por pregunta
(\cref{sec:disc-rq}) y la \cref{tab:teoria-experimento} cierran ese ciclo.

\paragraph{Archivos.} \texttt{report/main.tex}, \texttt{report/refs.bib} y
\texttt{report/sections/01-introduccion.tex}, \texttt{02-estado-del-arte.tex} y
\texttt{07-conclusiones.tex}.
\end{aporte}


\end{document}
